\documentclass[12pt,a4paper]{article}
\usepackage[utf8]{inputenc}
\usepackage{geometry}
\geometry{margin=1in}
\usepackage{setspace}
\onehalfspacing
\usepackage{parskip}
\usepackage{booktabs}
\usepackage{amsmath}
\usepackage{amssymb}
\usepackage{enumitem}
\usepackage{graphicx}
\usepackage{float}
\usepackage{algorithm}
\usepackage{algpseudocode}
\usepackage{array}
\usepackage{tabularx}
\usepackage{longtable}
\usepackage{listings}
\usepackage{xcolor}
\usepackage{hyperref}
\usepackage{url}

% Code listing style
\lstdefinelanguage{TypeScript}{
  keywords={import, export, from, const, let, var, function, async, await,
            return, if, else, switch, case, break, default, try, catch,
            throw, new, this, true, false, null, undefined, type, interface,
            Promise, string, number, boolean, void},
  keywordstyle=\color{blue}\bfseries,
  ndkeywords={doc, collection, addDoc, getDocs, getDoc, setDoc, updateDoc,
              query, where, orderBy, limit, deleteDoc, arrayUnion,
              NextResponse, NextRequest},
  ndkeywordstyle=\color{teal},
  sensitive=true,
  comment=[l]{//},
  morecomment=[s]{/*}{*/},
  commentstyle=\color{gray}\itshape,
  stringstyle=\color{orange},
  morestring=[b]",
  morestring=[b]',
  morestring=[b]`
}

\lstdefinelanguage{Python3}{
  keywords={import, from, def, class, return, if, else, elif, for, while,
            in, not, and, or, True, False, None, try, except, with, as,
            lambda, pass, raise, yield, async, await},
  keywordstyle=\color{blue}\bfseries,
  comment=[l]{\#},
  commentstyle=\color{gray}\itshape,
  stringstyle=\color{orange},
  morestring=[b]",
  morestring=[b]'
}

\lstset{
  basicstyle=\ttfamily\footnotesize,
  breaklines=true,
  frame=single,
  backgroundcolor=\color{gray!6},
  numbers=left,
  numberstyle=\tiny\color{gray},
  numbersep=6pt,
  xleftmargin=1.5em,
  tabsize=2,
  showstringspaces=false
}

\begin{document}

\begin{center}
    \textbf{\large CHAPTER 1}\\[0.4em]
    \textbf{\large INTRODUCTION}
\end{center}
\vspace{1.5em}

\noindent Canine blood transfusion is a life-saving intervention in veterinary emergency medicine, used in conditions including traumatic haemorrhage, immune-mediated haemolytic anaemia, surgical blood loss, and hemoperitoneum. Unlike human blood banking, which is supported by centralised national infrastructure, volunteer registries, and standardised logistics, canine transfusion services in India currently operate through independent veterinary clinics that maintain informal donor panels, rely on personal contacts for donor outreach, and lack any coordinated demand management system. The result is a fragmented network in which the availability of blood for a dog in a critical condition depends more on chance and personal connections than on structured, repeatable system design.

\noindent This project presents \textbf{K9Hope}, a web-based decision support platform developed to address the coordination gap in canine blood donation services across sparse veterinary networks in India. The platform was developed in the Department of Computer Science and Engineering, Rajalakshmi Institute of Technology, Chennai, in academic collaboration with Madras Veterinary College (MVC), Vepery. K9Hope integrates four functional components — a donor registry with automated Multi-Constraint Eligibility Filtering (MCEF), a Haversine-based Geospatial Matching Algorithm (GMA), a Natural Language Processing (NLP) document triage module, and a Zero-Inflated Conformal Prediction (ZICP) demand forecasting engine — within a unified, role-differentiated web application built on Next.js, Firebase, and Vercel. The platform was evaluated retrospectively on transfusion data from 12 veterinary clinics comprising 1,847 recorded transfusion events, and through simulation on 5,000 synthetic donor profiles against 500 emergency request scenarios.

\vspace{2em}
\noindent\textbf{1.1\quad PROBLEM STATEMENT}
\vspace{0.8em}

\noindent Veterinary transfusion centres in India currently operate without shared donor registries, structured eligibility screening workflows, demand forecasting infrastructure, or standardised inter-clinic communication protocols. When a dog requires a blood transfusion, clinical staff must manually contact personal networks, broadcast requests through informal messaging channels such as WhatsApp groups, and verify donor eligibility only after the donor physically arrives at the clinic. This process takes a median of 240 minutes and succeeds in identifying a suitable donor within the critical response window only 60\% of the time.

\noindent Several structural features make this problem difficult to address by simply scaling down human blood banking methods. Canine blood demand is primarily emergency-driven rather than scheduled. Across the 12 clinics observed in this study, 78\% of clinic-days recorded zero transfusion demand, with emergency spikes accounting for the majority of actual utilisation. Donor eligibility is biologically and ethically constrained: under the Department of Animal Husbandry and Dairying (DAHD) July 2025 Standard Operating Procedures, only dogs weighing at least 25 kg, aged between 1 and 8 years, with a Packed Cell Volume (PCV) of 35\% or above, screened negative for vector-borne pathogens, and whose owners provide informed consent are eligible to donate. Owner consent is revocable at any time, introducing scheduling uncertainty that has no equivalent in human volunteer donation systems.

\noindent Blood component shelf life compounds the urgency. Fresh whole blood must be used within 8 hours of collection if unrefrigerated. Stored whole blood remains viable for only 28--35 days, and platelets have a shelf life of approximately 5 days. Inventory decisions must therefore balance shortage risk against wastage — a trade-off that point-estimate forecasting methods such as ARIMA-family models cannot adequately support. ARIMA models detect weekly seasonality in demand but systematically fail to capture zero-inflated emergency spikes. Machine learning classifiers require sample sizes in the thousands for reliable training, whereas individual veterinary clinics record fewer than 200 transfusion events per year. Static safety-stock heuristics ignore clinic-specific demand patterns entirely, producing blood wastage rates of 15--25\% in practice.

\noindent In addition to these operational challenges, the DAHD 2025 Standard Operating Procedures introduced mandatory compliance requirements — including donor health documentation, donation interval enforcement, and DEA blood typing — that are operationally infeasible to track through paper records or informal communication channels. No existing digital platform in India provides an integrated solution for donor matching, eligibility compliance, and inventory management in the canine blood banking context.

\vspace{2em}
\noindent\textbf{1.2\quad SCOPE OF THE PROJECT}
\vspace{0.8em}

\noindent The K9Hope platform addresses the operational scope of canine blood donation coordination within the Indian veterinary healthcare context. The system is designed for deployment across networks of independent veterinary clinics and blood banks, with an initial implementation scope centred on the Chennai metropolitan region in partnership with Madras Veterinary College, Vepery.

\noindent The scope of this project encompasses the following functional boundaries:

\begin{itemize}[leftmargin=1.5em, itemsep=0.4em]
    \item \textbf{Donor registry and compliance:} The platform registers canine blood donors, enforces DAHD 2025 physical and haematological eligibility criteria at the point of registration, tracks donation history, and manages donor welfare through programmatic cooldown enforcement. The scope is limited to dogs as donor animals; extension to other species is identified as future work.
    \item \textbf{Blood request management and triage:} The platform accepts blood transfusion requests from authenticated veterinary clinics, processes uploaded veterinary recommendation letters through an NLP triage pipeline, and assigns urgency priority classifications to support coordinator decision-making. The scope does not include direct patient management or electronic health record integration.
    \item \textbf{Geospatial donor matching:} The platform identifies compatible, eligible donors within a configurable geographic proximity to the requesting clinic using Haversine distance computation. The matching output is a ranked shortlist presented to a human coordinator; automated donor contact is outside the system's scope.
    \item \textbf{Inventory and welfare management:} The platform tracks blood component shelf life across the clinic network, issues First-Expiry-First-Out (FEFO) alerts for near-expiry units, and enforces donor rest intervals. The scope does not include automated blood product procurement or integration with commercial blood bank supply chains.
    \item \textbf{Demand forecasting:} The ZICP module provides clinic-level daily blood demand prediction intervals to support weekly inventory ordering decisions. The scope is limited to short-term (one-week-ahead) probabilistic forecasting; multi-step forecasting and spatial demand correlation are identified as future extensions.
    \item \textbf{Evaluation:} The platform is evaluated through simulation on synthetic data and retrospective analysis on historical clinical records. Prospective clinical validation through operational deployment is outside the scope of this project.
\end{itemize}

\noindent The project does not replace or replicate the clinical functions of a veterinary professional. All decisions involving donor selection, transfusion approval, and patient management remain under the authority of licensed veterinarians. The platform functions exclusively as a decision support and coordination tool.

\vspace{2em}
\noindent\textbf{1.3\quad OBJECTIVES}
\vspace{0.8em}

\noindent The primary objective of this project is to design, implement, and evaluate a centralised decision support platform for canine blood donation coordination that reduces donor identification latency, enforces regulatory compliance, and supports inventory planning under conditions of extreme data sparsity. The specific objectives are as follows:

\begin{enumerate}[leftmargin=1.8em, itemsep=0.5em]
    \item To develop a donor registry module with a Multi-Constraint Eligibility Filter (MCEF) that enforces DAHD 2025 donor criteria — including weight, age, PCV thresholds, DEA blood type, pathogen screening status, and donation cooldown intervals — at the point of registration, eliminating the need for manual eligibility verification at the time of emergency.
    \item To implement a Haversine-based Geospatial Matching Algorithm (GMA) that identifies verified, eligible donors within a configurable proximity radius and ranks results by distance, minimising blood transport time and reducing the risk of component degradation during transit.
    \item To develop an NLP-based document triage module that scans uploaded veterinary recommendation letters for urgency indicators — including terms associated with trauma, severe anaemia, and acute haemorrhage — and assigns priority classifications to incoming blood requests to support clinical staff in allocating attention during concurrent emergencies.
    \item To propose and evaluate the Zero-Inflated Conformal Prediction (ZICP) method for blood demand forecasting, providing distribution-free prediction intervals with finite-sample coverage guarantees suited to the zero-inflated, low-volume nature of veterinary transfusion data, without requiring distributional assumptions or large training datasets.
    \item To implement a First-Expiry-First-Out (FEFO) inventory tracking module that monitors component shelf life across blood product types and issues depletion alerts to reduce wastage from expired units, improving inventory utilisation across the clinic network.
    \item To design the platform in accordance with a human-in-the-loop architecture, ensuring that algorithmic outputs serve as decision support rather than autonomous clinical action, and that all final transfusion decisions remain under the authority of licensed veterinary professionals.
    \item To evaluate system performance through simulation on 5,000 synthetic donor profiles and retrospective analysis on 1,847 transfusion records from 12 veterinary clinics, measuring time-to-match, donor identification success rate, eligibility filter accuracy, ZICP forecast coverage, inventory cost reduction, and blood wastage reduction as primary evaluation metrics.
\end{enumerate}

\vspace{2em}
\noindent\textbf{1.4\quad K9HOPE: PLATFORM OVERVIEW}
\vspace{0.8em}

\noindent K9Hope is a web-based decision support platform designed to coordinate canine blood donation services across sparse veterinary networks in India. The platform addresses a documented operational gap: when a dog requires an emergency blood transfusion, there is currently no structured digital system in India to connect the patient with a compatible, eligible, and geographically proximate donor. The manual search process relies on phone calls, personal contacts, and informal social media broadcasts — a process with a median duration of 240 minutes and a success rate of only 60\% within the critical response window.

\noindent The platform is structured around six integrated functional modules. The \textbf{Multi-Constraint Eligibility Filter (MCEF)} enforces DAHD 2025 donor criteria automatically at the point of registration, screening each donor against weight, age, PCV, DEA blood type, pathogen test status, and donation cooldown constraints. This ensures that all donors in the registry are pre-verified, eliminating the need for manual eligibility checks at the time of emergency. The \textbf{Geospatial Matching Algorithm (GMA)} uses the Haversine formula to compute great-circle distances between a requesting clinic and registered donors, returning a proximity-ranked shortlist of compatible, eligible candidates for coordinator review.

\noindent The \textbf{NLP Triage Module} processes uploaded veterinary recommendation letters using keyword extraction to identify urgency indicators — terms associated with trauma, acute haemorrhage, and severe anaemia — and assigns each incoming blood request a priority classification of Critical, High, or Standard. This allows coordinators to manage concurrent requests by clinical urgency rather than submission order. The \textbf{ZICP Demand Forecasting Engine} provides probabilistic prediction intervals for daily transfusion demand at the clinic level, using a two-component zero-inflated model that separately handles demand occurrence and magnitude. The method achieves 92\% empirical coverage with intervals that are 35\% tighter than bootstrap alternatives, enabling a 22\% reduction in inventory ordering costs.

\noindent The \textbf{FEFO Inventory Management Module} tracks the shelf life of blood components across the network and issues First-Expiry-First-Out alerts to coordinators when units approach their expiry threshold, reducing wastage from expired products. All modules operate within a \textbf{human-in-the-loop architecture}: every output produced by the system — whether a donor shortlist, a triage classification, or an inventory alert — is presented to a qualified veterinary professional for review before any action is taken. The platform does not send automated messages to donor owners, does not approve transfusions autonomously, and does not prescribe clinical decisions.

\noindent The system is implemented on a modern cloud-native stack. The frontend is built in Next.js with server-side rendering; the backend uses Firebase Firestore for real-time data synchronisation, Firebase Authentication for role-based access control across donor owner, veterinary clinic, and coordinator roles, and Uploadcare for secure storage of medical documents. The platform is deployed on Vercel. The entire system was developed as a final-year undergraduate research project at the Department of Computer Science and Engineering, Rajalakshmi Institute of Technology, Chennai, and has not been clinically deployed. All reported performance metrics are derived from retrospective analysis and simulation; the platform provides a validated architectural foundation for future operational partnership with veterinary institutions.

\vspace{2em}
\noindent\textbf{1.5\quad ORGANISATION OF THE REPORT}
\vspace{0.8em}

\noindent The remainder of this report is organised as follows:

\vspace{0.6em}
\noindent \textbf{Chapter 2 — Literature Survey} presents a structured review of prior work across four areas directly relevant to this project: veterinary transfusion systems and donor management; blood supply chain optimisation and demand forecasting under data sparsity; conformal prediction and distribution-free uncertainty quantification; and donor return behaviour classification using machine learning. The chapter identifies the specific gaps in the existing literature that this project targets and positions the K9Hope platform and ZICP method within the broader research landscape.

\vspace{0.6em}
\noindent \textbf{Chapter 3 — Methodology} describes the complete system design of the K9Hope platform. It presents the architecture diagram, explains the data flow across functional modules, and documents the UML diagrams — including use case, class, activity, sequence, data flow, and entity-relationship diagrams — that formalise the system's structure and operational flows. The mathematical formulations for the Haversine distance computation, the ZICP non-conformity scoring function, and the conformal calibration procedure are also presented here, along with descriptions of all algorithms implemented in the system.

\vspace{0.6em}
\noindent \textbf{Chapter 4 — Implementation} provides a detailed account of the implementation of each platform module. It specifies the software requirements, input and output data formats, hardware and software resource requirements, and the design decisions made during development. The chapter documents the implementation of the MCEF compliance rules, the GMA query execution logic, the NLP keyword extraction pipeline, the ZICP training and calibration procedure, and the FEFO alert generation logic, supported by relevant code excerpts.

\vspace{0.6em}
\noindent \textbf{Chapter 5 — Results and Discussion} presents the experimental evaluation of the platform across all primary metrics: time-to-match, donor identification success rate, eligibility filter accuracy, ZICP empirical coverage and interval width, inventory cost reduction, blood wastage reduction, and donor return classification performance. Results are presented with supporting figures and tables, discussed in terms of their operational significance, and qualified by the limitations of the simulation-based evaluation methodology.

\vspace{0.6em}
\noindent \textbf{Chapter 6 — Conclusion} summarises the contributions of the project and identifies the principal finding: that information fragmentation rather than algorithmic complexity is the dominant barrier to canine blood service coordination in India. Directions for future work are outlined, including prospective clinical deployment, PMS API integration, multi-step ZICP forecasting, and nationwide registry expansion.

\vspace{0.6em}
\noindent \textbf{Appendices} include sample source code for key system modules, platform screenshots illustrating the donor registration interface, coordinator dashboard, geospatial donor matching view, and inventory management panel, and the conference papers submitted from this project along with their acceptance certificates.

\newpage
\begin{center}
    \textbf{\large CHAPTER 2}\\[0.4em]
    \textbf{\large LITERATURE REVIEW}
\end{center}
\vspace{1.5em}

\noindent The coordination of blood supply chains in healthcare settings has been studied extensively in the context of human blood banking, where centralised infrastructure, large donor pools, and predictable demand patterns enable the application of robust optimisation and stochastic programming methods. However, canine transfusion medicine presents a fundamentally different operational environment — characterised by decentralised clinic-level records, biologically constrained donor eligibility, owner-mediated consent, and emergency-driven demand — that renders direct adaptation of human blood supply chain methods inappropriate. This chapter reviews the relevant body of literature across four areas: veterinary transfusion practice and donor management, blood supply chain optimisation and demand forecasting, uncertainty quantification and conformal prediction methods, and donor return behaviour prediction using machine learning.

\vspace{2em}
\noindent\textbf{2.1\quad VETERINARY TRANSFUSION MEDICINE AND DONOR MANAGEMENT}
\vspace{0.8em}

\noindent The clinical foundations of canine transfusion medicine are well established. Tocci \cite{tocci2009} provides a comprehensive overview of transfusion medicine principles in small animal practice, covering component selection, compatibility testing, and transfusion reactions. The 2020 AAHA Canine Blood Component Transfusion Guidelines by Brooks et al.\ \cite{brooks2020} represent the current clinical standard, specifying eligibility criteria for donor animals, blood typing protocols using the Dog Erythrocyte Antigen (DEA) system, component preparation procedures, and recommended storage conditions. These guidelines form the regulatory basis on which the eligibility constraints in this project's Multi-Constraint Eligibility Filter (MCEF) are defined.

\noindent Davidow et al.\ \cite{davidow2013} published the Association of Veterinary Hematology and Transfusion Medicine (AVHTM) Transfusion Reaction Small Animal Consensus Statement, which identifies the primary categories of transfusion reaction in dogs and establishes the clinical rationale for blood typing and cross-matching prior to every transfusion. The significance of DEA 1.1 typing in determining transfusion compatibility and reaction risk is specifically addressed by Sharma and Raj \cite{sharma2014}, whose work informs the blood type matching logic implemented in the platform's donor-recipient pairing module. McMichael et al.\ \cite{mcmichael2009} examined the effect of leukoreduction on transfusion-induced immunomodulation in dogs, establishing leukoreduction status as a clinically meaningful attribute for component tracking — a feature incorporated into the inventory management module of the present system.

\noindent The biological constraints on canine donors are more restrictive than their human equivalents. Lanevschi and Wardrop \cite{lanevschi2001} outline the principles of transfusion medicine in small animals, noting that the eligible donor pool is limited by weight thresholds, age ranges, behavioural suitability, and infectious disease screening requirements. In Indian conditions specifically, seasonal variation in canine vector-borne disease prevalence — including \textit{Babesia} and \textit{Ehrlichia} infections — complicates screening, as Kumar et al.\ \cite{kumar2022} document in their study of metropolitan India. Spada et al.\ \cite{spada2016} similarly demonstrate significant seroprevalence of vector-borne pathogens in donor populations, reinforcing the need for mandatory pathogen screening as a pre-donation requirement rather than a post-arrival check. The present system encodes all these constraints as automated pre-filters applied at donor registration, eliminating the delay caused by eligibility assessment at the point of emergency.

\noindent A 2025 multicenter study by Dini et al.\ \cite{dini2025} examined physiological and behavioural changes in canine blood donors across donation events, providing empirical evidence that donor welfare is measurably affected by donation frequency and handling conditions. This finding underpins the donation interval enforcement logic in the platform, which enforces the minimum 30-day cooldown mandated by the DAHD 2025 Standard Operating Procedures \cite{dahd2025, pib2025} and tracks each donor's next eligible donation date to prevent over-donation.

\vspace{2em}
\noindent\textbf{2.2\quad BLOOD SUPPLY CHAIN COORDINATION AND DEMAND FORECASTING}
\vspace{0.8em}

\noindent The broader literature on blood supply chain optimisation provides methodological precedents, though most existing work assumes conditions — large historical datasets, centralised data access, and relatively stable demand — that do not hold in sparse veterinary networks. Zahiri and Pishvaee \cite{zahiri2018} formulate a blood supply chain network design problem under uncertainty, incorporating blood group compatibility constraints into a stochastic programming model. Butt et al.\ \cite{butt2018} address blood transfusion demand and allocation optimisation in emergency settings using operations research methods. While these frameworks are not directly transferable, they provide the conceptual basis for treating canine blood inventory as a perishable supply chain optimisation problem with shortage and wastage cost components.

\noindent Dehghani et al.\ \cite{dehghani2021} propose proactive transshipment strategies for the blood supply chain using stochastic programming, demonstrating that coordinated inter-facility inventory sharing reduces both shortage and wastage rates. This finding motivates the cross-clinic inventory visibility feature in the coordination interface of the present system, which surfaces unit availability across the network rather than restricting visibility to individual clinic records. Khakshouri Fariman et al.\ \cite{khakshouri2024} extend multi-objective blood supply chain optimisation under scenario-based uncertainty, further establishing the theoretical case for uncertainty-aware inventory planning rather than point-estimate-based safety stock rules.

\noindent For intermittent demand, Croston's method and its variants (SBA, TSB) are the standard baseline, separating the forecasting of inter-demand intervals from demand sizes \cite{croston1972}. However, as Liu \cite{liu2020} demonstrates in the context of medical consumables with short life cycles during the COVID-19 epidemic, intermittent demand forecasting in healthcare settings is complicated by structural non-stationarity and zero-inflation that Croston-based methods do not explicitly model. Dinamarca et al.\ \cite{dinamarca2025} propose SARIMAX models with zero-inflated skew-normal errors as an extension for zero-inflated time series, though these parametric approaches require distributional assumptions that may not hold with fewer than 200 observations per clinic per year — the typical data volume in the veterinary setting studied here.

\noindent Mutimbanyoka and Ndlovu \cite{mutimbanyoka2025} apply predictive modelling to blood demand forecasting in a low-resource context, finding that model performance degrades substantially under data sparsity and that uncertainty quantification is more operationally valuable than point accuracy in such settings. This conclusion directly motivates the adoption of conformal prediction as the forecasting framework in the present work, in preference to point-estimate models that produce no actionable information about forecast uncertainty.

\vspace{2em}
\noindent\textbf{2.3\quad CONFORMAL PREDICTION AND UNCERTAINTY QUANTIFICATION}
\vspace{0.8em}

\noindent Conformal prediction, introduced formally by Vovk, Gammerman, and Shafer \cite{vovk2005}, provides a framework for constructing prediction intervals with guaranteed finite-sample coverage under the assumption of data exchangeability, without requiring distributional assumptions about the underlying process. Angelopoulos and Bates \cite{angelopoulos2021} provide a comprehensive introduction to conformal prediction and distribution-free uncertainty quantification, establishing its applicability across regression, classification, and time series settings. Romano et al.\ \cite{romano2019} extend the framework to conformalized quantile regression, enabling asymmetric prediction intervals that better capture skewed distributions — a relevant property for zero-inflated demand data where the lower bound of any valid interval is zero on the majority of days.

\noindent Stankeviciute et al.\ \cite{stankeviciute2021} develop conformal prediction methods for time series forecasting, addressing the exchangeability challenge posed by temporal autocorrelation. Xu and Xie \cite{xu2023} extend this to formal guarantees for time series conformal prediction under mild distributional conditions. Zaffran et al.\ \cite{zaffran2022} propose adaptive conformal predictions for time series that update the calibration threshold in response to distribution shift — a mechanism that is particularly relevant to veterinary demand data, where epidemic events such as canine parvovirus outbreaks can produce sustained deviations from historical patterns.

\noindent Lei et al.\ \cite{lei2018} establish the theoretical foundations for distribution-free predictive inference in regression settings, demonstrating that conformal methods achieve coverage guarantees that hold for any finite sample size — a critical property when training data consists of fewer than 100 observations per clinic. Vazquez and Facelli \cite{vazquez2022} review the application of conformal prediction in clinical medical sciences, demonstrating coverage reliability in settings with small and imbalanced datasets. These properties collectively justify the selection of conformal prediction as the methodological foundation for the ZICP method proposed in this project, in preference to Bayesian neural networks or bootstrap methods that require either strong prior specification or large sample sizes to produce reliable intervals.

\vspace{2em}
\noindent\textbf{2.4\quad DONOR RETURN BEHAVIOUR AND CLASSIFICATION}
\vspace{0.8em}

\noindent Predicting whether a blood donor will return for a subsequent donation is an established problem in human transfusion medicine. Yeh et al.\ \cite{yeh2025} apply machine learning methods to predict blood donor return behaviour in human settings, evaluating multiple classifiers and finding that ensemble methods outperform logistic regression when sufficient training data is available. Their analysis of feature importance — identifying recency of last donation, frequency of past donations, and cumulative volume as primary predictors — provides the feature engineering rationale adopted in the donor return classification module of the present system.

\noindent However, the canine setting introduces complications absent from human donor prediction. In the present study, donor return is mediated by owner decision rather than donor self-motivation, meaning that owner-level variables — proximity to clinic, socioeconomic factors, prior experience with the donation process — likely dominate the return decision but are not captured in administrative records. With 176 donors in the retrospective dataset, classification models face severe overfitting risk; the ANN model achieves 78.4\% accuracy compared to 76.5\% for Random Forest, a difference of 1.9 percentage points that is statistically significant at $p < 0.05$ but operationally modest given the sample size. These results are reported transparently as establishing a methodological baseline rather than as a deployable prediction tool, and are consistent with findings in the sparse healthcare prediction literature that recommend interpretable models such as logistic regression and Random Forest in low-data settings over complex architectures that are difficult to validate \cite{chawla2002}.

\vspace{2em}
\noindent\textbf{2.5\quad IDENTIFIED GAPS AND POSITIONING OF THE PRESENT WORK}
\vspace{0.8em}

\noindent The reviewed literature collectively establishes that: (i) canine transfusion medicine has well-defined clinical standards but no coordinated digital infrastructure for donor-recipient matching in India; (ii) blood supply chain optimisation methods exist but assume data volumes and centralisation not present in veterinary settings; (iii) conformal prediction provides coverage guarantees suitable for sparse, zero-inflated demand but has not previously been applied to veterinary blood inventory; and (iv) donor return prediction in the canine context is confounded by owner-mediated consent and the absence of owner-level variables in clinical records. The present project addresses gaps (i) and (iii) through the K9Hope platform architecture and the ZICP forecasting method, while acknowledging the limitations identified in (ii) and (iv) as directions for future work requiring larger datasets and operational deployment partnerships.

\newpage
\begin{center}
    \textbf{\large CHAPTER 3}\\[0.4em]
    \textbf{\large METHODOLOGY}
\end{center}
\vspace{1.5em}

\noindent This chapter presents the complete system design of the K9Hope platform. The methodology is structured around six interdependent functional modules, each addressing a specific coordination bottleneck identified in Chapter 1. The design philosophy throughout is conservative and human-centred: every algorithmic output is a recommendation surfaced for review by a qualified veterinary professional, not an autonomous action. The chapter proceeds from the high-level architecture diagram through module-level descriptions, data flow, entity-relationship modelling, UML diagrams, and the formal algorithmic specifications underlying the system's two primary analytical components — the Haversine-based Geospatial Matching Algorithm and the Zero-Inflated Conformal Prediction demand forecasting method.

\vspace{2em}
\noindent\textbf{3.1\quad ARCHITECTURE DIAGRAM}
\vspace{0.8em}

\noindent The K9Hope platform is organised as a layered, role-differentiated web application with three user-facing roles — Donor Owner, Veterinary Clinic, and Coordinator — each interacting with the system through separate authenticated dashboards. The backend is implemented as a set of Firebase Cloud Functions that expose API endpoints consumed by the Next.js frontend; persistent data is stored in Firebase Firestore in real time, and medical documents are stored in Uploadcare with access tokens bound to authenticated user sessions.

\noindent The system architecture is illustrated in Figure 3.1. Data flows from donor registration and blood request submission through a pipeline of automated processing stages — eligibility filtering, NLP triage, geospatial matching, and inventory management — before reaching the Coordinator dashboard, where every match, priority classification, and alert is reviewed before any action is taken. The ZICP forecasting module operates as a parallel analytical service that consumes historical demand records from Firestore and returns prediction intervals to the inventory management dashboard on a scheduled basis.

\begin{figure}[H]
    \centering
    \includegraphics[width=\linewidth]{fig_3_1_architecture}
    \caption{\textbf{Fig. 3.1 K9Hope System Architecture — Layered Role-Differentiated Platform}}
\end{figure}

\vspace{2em}
\noindent\textbf{3.2\quad SYSTEM DESIGN OVERVIEW}
\vspace{0.8em}

\noindent The K9Hope platform is designed around a central principle: the primary barrier to canine blood service coordination is information fragmentation, not algorithmic complexity. The architecture therefore prioritises data integration, compliance automation, and visibility over predictive sophistication. The system maintains a single source of truth — a Firebase Firestore database — that is accessible in real time to all authenticated users, replacing the fragmented landscape of phone records, WhatsApp broadcast groups, and paper donor registers that currently characterise Indian veterinary blood banking.

\noindent The platform enforces a strict separation between algorithmic processing and clinical decision-making. All computationally derived outputs — donor shortlists, urgency classifications, demand intervals, expiry alerts — are flagged as recommendations and presented to a Coordinator for explicit approval before any downstream action is triggered. No automated message is sent to a donor owner without Coordinator approval. No transfusion is initiated by the system. This design is directly informed by the ethical constraints of canine donation, where owner consent is revocable at any time and cannot be assumed, and by the epistemic limitations of the classification and forecasting models, which are trained on sparse data and carry non-trivial uncertainty.

\noindent The technology stack is as follows. The frontend is built in Next.js 14 with TypeScript and server-side rendering, styled with Tailwind CSS and shadcn/ui. Firebase Firestore provides the real-time NoSQL database; Firebase Authentication handles role-based access control through JWT token verification on all API routes. Uploadcare provides secure, CDN-backed storage for medical documents with time-limited access tokens. The ZICP forecasting module is implemented in Python 3.10 using LightGBM 4.0 and scikit-learn 1.3 and is invoked as a scheduled Cloud Function. The platform is deployed on Vercel with edge caching enabled for authenticated dashboard routes.

\vspace{2em}
\noindent\textbf{3.3\quad MODULES DESCRIPTION}
\vspace{0.8em}

\noindent The K9Hope platform is implemented as six functional modules. Each module is described below in terms of its purpose, inputs, processing logic, and outputs.

\vspace{1em}
\noindent\textbf{Module 1: Donor Registry and Multi-Constraint Eligibility Filter (MCEF)}
\vspace{0.5em}

\noindent The donor registry module is the entry point for all donor-side activity. A pet owner creates an account, selects the Donor Owner role, and submits a registration form providing the dog's name, breed, date of birth, body weight (kg), sex, DEA blood type, most recent PCV value (\%), pathogen screening status (positive/negative for Babesia, Ehrlichia, Leishmania, Heartworm), and vaccination history. The MCEF evaluates each submitted profile against DAHD 2025 eligibility criteria at the point of form submission, before any record is written to the database.

\noindent The MCEF enforces the following constraints simultaneously:

\begin{itemize}[leftmargin=1.5em, itemsep=0.3em]
    \item Weight $\geq$ 25 kg
    \item Age: 1 year $\leq$ age $\leq$ 8 years (computed from date of birth)
    \item PCV $\geq$ 35\%
    \item Pathogen screening: negative for all four vector-borne pathogens
    \item Donation cooldown: most recent donation date $\geq$ 30 days before current date
    \item Owner consent: digitally confirmed with timestamp at submission
\end{itemize}

\noindent If any constraint is violated, the form returns a specific field-level validation error identifying the failed criterion. The profile is not written to the registry. If all constraints pass, the donor profile is written to Firestore with a verified status flag, and the record becomes available for geospatial matching. Cooldown enforcement is re-evaluated dynamically at match time: a donor who was eligible at registration may fall into a cooldown period following a subsequent donation.

\vspace{1em}
\noindent\textbf{Module 2: Blood Request Submission and NLP Triage}
\vspace{0.5em}

\noindent Veterinary clinics submit blood requests through an authenticated clinic dashboard. The request form captures patient details (breed, weight, age, DEA blood type required), clinical indication (free text), required blood component (whole blood, packed RBCs, plasma, platelets), and an uploaded veterinary recommendation letter (PDF or image). On submission, the NLP triage module processes the uploaded document using Optical Character Recognition (OCR) to extract plain text, followed by keyword matching against a predefined urgency lexicon.

\noindent The urgency lexicon is organised into three tiers:

\begin{itemize}[leftmargin=1.5em, itemsep=0.3em]
    \item \textbf{Critical:} trauma, hemoperitoneum, acute haemorrhage, haemodynamic instability, GCS $\leq$ 8, PCV $<$ 15\%, Hb $<$ 5 g/dL, collapse, shock
    \item \textbf{High:} severe anaemia, immune-mediated haemolytic anaemia, IMHA, surgical blood loss, transfusion-dependent, PCV 15--25\%
    \item \textbf{Standard:} elective procedure, pre-operative, chronic anaemia, monitoring
\end{itemize}

\noindent The highest-matching tier determines the request's priority classification. The classified request, along with extracted urgency terms highlighted for coordinator review, is surfaced on the Coordinator dashboard. A Coordinator reviews the classification and the original document before approving the request for matching. The NLP module does not make clinical decisions; its output is a prioritised queue that allows coordinators to allocate attention to the most urgent cases during concurrent emergencies.

\vspace{1em}
\noindent\textbf{Module 3: Haversine-based Geospatial Matching Algorithm (GMA)}
\vspace{0.5em}

\noindent Upon coordinator approval of a blood request, the GMA queries the donor registry for all profiles that satisfy the following simultaneous constraints: (1) DEA blood type compatible with the requested type, (2) MCEF eligibility status verified and cooldown not active, (3) geographic distance from the requesting clinic within a configurable radius (default: 10 km). Distance is computed using the Haversine formula applied to the GPS coordinates stored at donor registration and the clinic's registered coordinates.

\noindent The Haversine formula computes the great-circle distance $d$ between two points $(\phi_1, \lambda_1)$ and $(\phi_2, \lambda_2)$ on a sphere of radius $R$ as:

\[
a = \sin^2\!\left(\frac{\phi_2 - \phi_1}{2}\right) + \cos\phi_1 \cdot \cos\phi_2 \cdot \sin^2\!\left(\frac{\lambda_2 - \lambda_1}{2}\right)
\]

\[
d = 2R \cdot \arcsin\!\left(\sqrt{a}\right)
\]

\noindent where $\phi$ denotes latitude in radians, $\lambda$ denotes longitude in radians, and $R = 6{,}371$ km is the mean Earth radius. The resulting shortlist is sorted by ascending distance and presented to the Coordinator. The Coordinator selects donors for outreach from this list; the system then opens a secure, authenticated messaging channel between the Coordinator and the selected donor owner's account. No automated contact message is sent.

\vspace{1em}
\noindent\textbf{Module 4: ZICP Demand Forecasting Engine}
\vspace{0.5em}

\noindent The ZICP module provides clinic-level daily blood demand prediction intervals to support weekly inventory ordering decisions. It is implemented as a Python service that runs on a scheduled basis, consuming the previous 90 days of transfusion records from Firestore and writing the output prediction intervals back to a Firestore forecasting collection accessible to the Coordinator dashboard.

\noindent The module models demand using a zero-inflated distribution with two separately trained components. The occurrence model is a LightGBM binary classifier that predicts the probability of any demand occurring on a given day. The size model is a LightGBM log-normal regressor trained on positive-demand records only, estimating the mean and standard deviation of log demand. Prediction intervals are calibrated using held-out data through the conformal calibration procedure described in Section 3.8. The full mathematical formulation is presented in Section 3.8.

\vspace{1em}
\noindent\textbf{Module 5: FEFO Inventory Management}
\vspace{0.5em}

\noindent The inventory management module maintains a unit-level record for each blood component in the clinic network. Each unit record stores: component type (whole blood, packed RBCs, FFP, platelets, cryoprecipitate), DEA blood type, collection date, leukoreduction status, storage location, and assigned expiry date computed from component type. The module computes remaining shelf life for each unit on a daily basis and generates First-Expiry-First-Out (FEFO) alerts for units with fewer than 3 days of remaining shelf life. Alerts are surfaced on the Coordinator dashboard in order of ascending remaining shelf life. When a Coordinator allocates a unit to a patient, the unit record is marked as dispensed, removing it from the active inventory.

\vspace{1em}
\noindent\textbf{Module 6: Coordinator Dashboard and Communication Layer}
\vspace{0.5em}

\noindent The Coordinator dashboard aggregates outputs from all five upstream modules into a single role-specific interface. The dashboard presents: the NLP-triaged blood request queue, ranked by urgency classification; the GMA donor shortlists for approved requests; FEFO inventory alerts; ZICP prediction interval charts for the current week; and a secure messaging panel for approved donor outreach. All interactions on the Coordinator dashboard are logged with timestamps to Firestore for auditability. The messaging layer uses Firebase Realtime Database for low-latency delivery of messages between Coordinators and donor owners; message threads are request-scoped and cannot be initiated by either party outside of an approved request context.

\vspace{2em}
\noindent\textbf{3.4\quad DATA FLOW DIAGRAM}
\vspace{0.8em}

\noindent The Data Flow Diagrams (DFDs) presented below model the movement of data through the K9Hope system at two levels of abstraction.

\noindent\textbf{Level 0 DFD (Context Diagram):} Figure 3.2 shows the system as a single process receiving inputs from three external entities — Donor Owner, Veterinary Clinic, and Coordinator — and producing outputs back to each. The Donor Owner provides registration data and receives eligibility confirmation or rejection. The Veterinary Clinic provides blood requests and receives match shortlists and inventory status. The Coordinator receives triage queues, donor shortlists, FEFO alerts, and forecast intervals, and produces approval decisions and outreach initiations.

\begin{figure}[H]
    \centering
    \includegraphics[width=0.85\linewidth]{fig_3_2_dfd_level0}
    \caption{\textbf{Fig. 3.2 Level 0 Data Flow Diagram (Context Diagram)}}
\end{figure}

\noindent\textbf{Level 1 DFD:} Figure 3.3 decomposes the system into five internal processes corresponding to the five automated modules. Data stores (Donor Registry, Request Queue, Inventory Records, Demand History) are shown as persistent repositories that mediate data exchange between processes. The MCEF process reads from and writes to the Donor Registry data store. The NLP Triage process reads uploaded documents and writes classified requests to the Request Queue. The GMA process reads from both the Donor Registry and Request Queue and writes ranked shortlists to the Coordinator interface. The ZICP process reads from Demand History and writes forecast intervals to the Inventory module. The FEFO process reads from and writes alerts to Inventory Records.

\begin{figure}[H]
    \centering
    \includegraphics[width=\linewidth]{fig_3_3_dfd_level1}
    \caption{\textbf{Fig. 3.3 Level 1 Data Flow Diagram}}
\end{figure}

\vspace{2em}
\noindent\textbf{3.5\quad E-R DIAGRAM}
\vspace{0.8em}

\noindent The Entity-Relationship (E-R) diagram in Figure 3.4 models the persistent data entities of the K9Hope platform and the relationships between them. The primary entities are as follows:

\begin{itemize}[leftmargin=1.5em, itemsep=0.4em]
    \item \textbf{User:} Stores authentication identity, role (DonorOwner / VetClinic / Coordinator), and contact details. Primary key: \texttt{userID}.
    \item \textbf{DonorProfile:} Stores canine biological attributes — name, breed, DOB, weight, sex, DEA type, PCV, pathogen screen results, consent timestamp, MCEF eligibility status, and last donation date. Foreign key: \texttt{ownerID} references \texttt{User.userID}.
    \item \textbf{BloodRequest:} Stores request details — requesting clinic ID, patient DEA type required, component type, clinical indication, uploaded document URL, NLP urgency classification, and Coordinator approval status. Foreign key: \texttt{clinicID} references \texttt{User.userID}.
    \item \textbf{DonationRecord:} Stores completed donation events — donor profile ID, matched request ID, donation date, component type, volume collected (mL), and collection site. Foreign keys reference both \texttt{DonorProfile} and \texttt{BloodRequest}.
    \item \textbf{InventoryUnit:} Stores individual blood unit records — component type, DEA type, collection date, expiry date, leukoreduction status, storage location, and dispensed status. Foreign key: \texttt{donationID} references \texttt{DonationRecord}.
    \item \textbf{DemandRecord:} Stores daily aggregate transfusion demand per clinic per blood type, used as input to the ZICP forecasting module. Foreign key: \texttt{clinicID} references \texttt{User.userID}.
    \item \textbf{MessageThread:} Stores secure communication threads between Coordinators and Donor Owners, scoped to a specific \texttt{BloodRequest}.
\end{itemize}

\begin{figure}[H]
    \centering
    \includegraphics[width=\linewidth]{fig_3_4_er_diagram}
    \caption{\textbf{Fig. 3.4 Entity-Relationship Diagram — K9Hope Data Model}}
\end{figure}

\vspace{2em}
\noindent\textbf{3.6\quad USE CASE AND CLASS DIAGRAM}
\vspace{0.8em}

\noindent\textbf{Use Case Diagram}
\vspace{0.5em}

\noindent Figure 3.5 presents the Use Case Diagram for the K9Hope platform. The system boundary encompasses five primary use cases: Register Donor, Submit Blood Request, Review and Approve Request, Initiate Donor Outreach, and Manage Inventory. Three actors interact with the system: Donor Owner, Veterinary Clinic, and Coordinator. The Coordinator extends the Review and Approve Request use case with two sub-cases: Trigger Geospatial Match and Approve Donor Contact. The NLP Triage and MCEF Filter processes are shown as system-internal actors (secondary actors) that participate in the Submit Blood Request and Register Donor use cases respectively.

\begin{figure}[H]
    \centering
    \includegraphics[width=\linewidth]{fig_3_5_use_case}
    \caption{\textbf{Fig. 3.5 Use Case Diagram — K9Hope Platform}}
\end{figure}

\noindent\textbf{Class Diagram}
\vspace{0.5em}

\noindent Figure 3.6 presents the Class Diagram for the system's backend service layer. The primary classes are \texttt{DonorRegistry}, \texttt{MCEFFilter}, \texttt{BloodRequestQueue}, \texttt{NLPTriageModule}, \texttt{GeospatialMatcher}, \texttt{ZICPForecaster}, \texttt{InventoryManager}, and \texttt{CoordinatorService}. \texttt{MCEFFilter} has a dependency relationship with \texttt{DonorRegistry}: it validates incoming profiles and calls \texttt{DonorRegistry.writeVerifiedProfile()} on pass. \texttt{GeospatialMatcher} has an association with both \texttt{DonorRegistry} (to query eligible donors) and \texttt{BloodRequestQueue} (to read approved requests). \texttt{ZICPForecaster} composes two sub-classes — \texttt{OccurrenceModel} and \texttt{SizeModel} — and depends on \texttt{ConformalCalibrator} for interval construction. \texttt{InventoryManager} subscribes to alerts generated by \texttt{FEFOAlertService}, which is a component of \texttt{InventoryManager}.

\begin{figure}[H]
    \centering
    \includegraphics[width=\linewidth]{fig_3_6_class_diagram}
    \caption{\textbf{Fig. 3.6 Class Diagram — Backend Service Layer}}
\end{figure}

\vspace{2em}
\noindent\textbf{3.7\quad SEQUENCE DIAGRAM}
\vspace{0.8em}

\noindent Two sequence diagrams are presented, covering the two most critical operational flows in the system.

\noindent\textbf{Sequence Diagram 1: Emergency Blood Request to Donor Shortlist}
\vspace{0.5em}

\noindent Figure 3.7 models the message sequence from blood request submission by a Veterinary Clinic to the delivery of a ranked donor shortlist to the Coordinator dashboard. The sequence is as follows: (1) VetClinic submits request form with uploaded recommendation letter to the API; (2) API invokes NLPTriageModule.classify(document), which returns an urgency tier and extracted keywords; (3) API writes the classified request to Firestore BloodRequest collection; (4) Coordinator receives a real-time notification via Firestore listener; (5) Coordinator reviews and approves the request; (6) Approval triggers GeospatialMatcher.match(requestID), which queries DonorRegistry for compatible, eligible donors within radius; (7) GMA returns ranked shortlist; (8) Coordinator selects donors and initiates outreach, which opens a MessageThread in Firebase Realtime Database.

\begin{figure}[H]
    \centering
    \includegraphics[width=\linewidth]{fig_3_7_sequence_request}
    \caption{\textbf{Fig. 3.7 Sequence Diagram — Emergency Blood Request to Donor Shortlist}}
\end{figure}

\noindent\textbf{Sequence Diagram 2: Donor Registration with MCEF Validation}
\vspace{0.5em}

\noindent Figure 3.8 models the message sequence from donor registration submission to registry confirmation or rejection. The sequence is as follows: (1) DonorOwner submits registration form to API; (2) API invokes MCEFFilter.validate(profile), which evaluates all six eligibility constraints in sequence; (3) If any constraint fails, MCEFFilter returns a structured error identifying the failing field, and the API responds with a 422 validation error — no record is written; (4) If all constraints pass, MCEFFilter returns a verified status, and the API writes the profile to Firestore DonorRegistry with \texttt{eligibility: "verified"} and \texttt{cooldownUntil: null}; (5) DonorOwner receives confirmation and access to the donor dashboard.

\begin{figure}[H]
    \centering
    \includegraphics[width=0.85\linewidth]{fig_3_8_sequence_registration}
    \caption{\textbf{Fig. 3.8 Sequence Diagram — Donor Registration with MCEF Validation}}
\end{figure}

\vspace{2em}
\noindent\textbf{3.8\quad ALGORITHMS}
\vspace{0.8em}

\noindent This section presents the formal algorithmic specifications for the two primary analytical components of the K9Hope platform: the Haversine-based Geospatial Matching Algorithm and the Zero-Inflated Conformal Prediction method.

\noindent\textbf{Algorithm 1: Multi-Constraint Eligibility Filter (MCEF)}
\vspace{0.5em}

\noindent The MCEF evaluates donor profiles against DAHD 2025 criteria using a short-circuit logical chain. Constraints are evaluated in order of computational cost, with Boolean checks executed before database lookups.

\begin{algorithm}[H]
\caption{Multi-Constraint Eligibility Filter (MCEF)}
\begin{algorithmic}[1]
\Require Donor profile $P = \{w, a, \text{pcv}, \text{path}, \text{lastDon}, \text{consent}\}$
\Ensure Eligibility decision $E \in \{\textit{PASS}, \textit{FAIL}(\text{reason})\}$
\If{$w < 25$} \Return \textit{FAIL}(``Weight below 25 kg'') \EndIf
\If{$a < 1$ \textbf{or} $a > 8$} \Return \textit{FAIL}(``Age out of range 1--8 years'') \EndIf
\If{$\text{pcv} < 35$} \Return \textit{FAIL}(``PCV below 35\%'') \EndIf
\If{$\text{path} \neq \textit{negative\_all}$} \Return \textit{FAIL}(``Pathogen screen positive'') \EndIf
\If{$\text{consent} = \textit{false}$} \Return \textit{FAIL}(``Owner consent not recorded'') \EndIf
\If{$(\text{today} - \text{lastDon}) < 30$} \Return \textit{FAIL}(``Within 30-day cooldown'') \EndIf
\State \Return \textit{PASS}
\end{algorithmic}
\end{algorithm}

\vspace{1.2em}
\noindent\textbf{Algorithm 2: Haversine-based Geospatial Matching Algorithm (GMA)}
\vspace{0.5em}

\noindent Let $R$ denote the blood request with required DEA type $b_R$ and clinic coordinates $(\phi_c, \lambda_c)$. Let $\mathcal{D}$ denote the set of all donors in the registry. The GMA returns a ranked shortlist $\mathcal{S}$ of compatible, eligible donors within radius $r$ (default 10 km), sorted by ascending Haversine distance.

\begin{algorithm}[H]
\caption{Haversine-based Geospatial Matching Algorithm (GMA)}
\begin{algorithmic}[1]
\Require Request $R$, Donor set $\mathcal{D}$, Radius $r$, Earth radius $R_E = 6371$ km
\Ensure Ranked shortlist $\mathcal{S}$
\State $\mathcal{S} \leftarrow \emptyset$
\For{each donor $d \in \mathcal{D}$}
    \If{$d.\text{DEA} \neq b_R$} \textbf{continue} \EndIf
    \If{$d.\text{eligibility} \neq \textit{verified}$} \textbf{continue} \EndIf
    \If{$(\text{today} - d.\text{lastDon}) < 30$} \textbf{continue} \EndIf
    \State $\Delta\phi \leftarrow d.\phi - \phi_c$; \quad $\Delta\lambda \leftarrow d.\lambda - \lambda_c$
    \State $a \leftarrow \sin^2\!\left(\frac{\Delta\phi}{2}\right) + \cos(\phi_c) \cdot \cos(d.\phi) \cdot \sin^2\!\left(\frac{\Delta\lambda}{2}\right)$
    \State $\text{dist} \leftarrow 2 R_E \cdot \arcsin\!\left(\sqrt{a}\right)$
    \If{$\text{dist} \leq r$}
        \State $\mathcal{S} \leftarrow \mathcal{S} \cup \{(d,\, \text{dist})\}$
    \EndIf
\EndFor
\State \Return $\mathcal{S}$ sorted by ascending dist
\end{algorithmic}
\end{algorithm}

\vspace{1.2em}
\noindent\textbf{Algorithm 3: Zero-Inflated Conformal Prediction (ZICP) — Training and Calibration}
\vspace{0.5em}

\noindent Let $D_{c,b,t} \in \mathbb{N}_0$ denote the transfusion demand at clinic $c$, blood type $b$, day $t$. Demand follows a zero-inflated distribution:

\[
D = Z \cdot Y, \quad Z \sim \text{Bernoulli}(\pi), \quad Y \mid Y > 0 \sim \text{LogNormal}(\mu, \sigma^2)
\]

\noindent where $Z$ is the occurrence indicator and $Y$ is the demand magnitude given occurrence. The feature vector is $X_{c,b,t} = (\text{day-of-week},\ \text{month},\ \text{clinicID},\ \text{bloodType},\ \bar{D}_{7\text{-day}})$.

\noindent The non-conformity score is defined as:

\[
S(X, D) = \begin{cases}
-\log \hat{\pi}(X) & \text{if } D = 0 \\
-\log(1 - \hat{\pi}(X)) + \dfrac{(\log D - \hat{\mu}(X))^2}{2\hat{\sigma}^2(X)} & \text{if } D > 0
\end{cases}
\]

\noindent where $\hat{\pi}(X)$ is the predicted probability of zero demand, and $\hat{\mu}(X)$, $\hat{\sigma}(X)$ are the predicted mean and standard deviation of log positive demand. For $D = 0$, the score is low when the model correctly predicts high zero-demand probability. For $D > 0$, the score is low when both occurrence and magnitude are accurately predicted.

\noindent The conformal calibration procedure selects the empirical quantile of calibration scores as the threshold for interval construction:

\[
q = \text{Quantile}\!\left(\{s_i\}_{i=1}^{n_\text{cal}},\ \frac{\lceil(n_\text{cal}+1)(1-\alpha)\rceil}{n_\text{cal}}\right)
\]

\noindent The prediction interval for a new observation $X_{\text{new}}$ is:

\[
C(X_{\text{new}}) = \{d \in \mathbb{N}_0 : S(X_{\text{new}}, d) \leq q\}
\]

\noindent Under exchangeability of calibration and test data, the coverage guarantee holds:

\[
\mathbb{P}\!\left(D_{\text{new}} \in C(X_{\text{new}})\right) \geq 1 - \alpha \quad \text{for all sample sizes } n
\]

\begin{algorithm}[H]
\caption{ZICP Training, Calibration, and Interval Construction}
\begin{algorithmic}[1]
\Require Training data $\mathcal{D}_\text{train}$, Calibration data $\mathcal{D}_\text{cal}$, Significance level $\alpha$
\Ensure Calibrated quantile $q$, Trained models $\hat{\pi}$, $\hat{\mu}$, $\hat{\sigma}$
\State Train LightGBM classifier $f$ on $\mathcal{D}_\text{train}$ to predict $\mathbb{1}[D > 0]$
\State $\hat{\pi}(X) \leftarrow f(X)$ \Comment{Occurrence probability model}
\State $\mathcal{D}_\text{pos} \leftarrow \{(X, D) \in \mathcal{D}_\text{train} : D > 0\}$
\State Train LightGBM regressor $g$ on $\mathcal{D}_\text{pos}$ to predict $\log D$
\State Train LightGBM regressor $h$ on $\mathcal{D}_\text{pos}$ to predict $|\log D - g(X)|$
\State $\hat{\mu}(X) \leftarrow g(X)$; \quad $\hat{\sigma}(X) \leftarrow h(X)$
\For{each $(X_i, D_i) \in \mathcal{D}_\text{cal}$}
    \State Compute $s_i \leftarrow S(X_i, D_i;\, \hat{\pi}, \hat{\mu}, \hat{\sigma})$ using Equation above
\EndFor
\State $q \leftarrow \left\lceil(n_\text{cal}+1)(1-\alpha)\right\rceil / n_\text{cal}$-th order statistic of $\{s_i\}$
\State \Return $q,\, \hat{\pi},\, \hat{\mu},\, \hat{\sigma}$
\Statex
\textbf{At inference for new observation $X_\text{new}$:}
\State Find $D_\text{max}$ such that $S(X_\text{new}, D_\text{max}) \leq q$ \Comment{Binary search on $D$}
\State \Return interval $[L, U] = [0,\, D_\text{max}]$
\State Order quantity: $q_{c,b,t} \leftarrow \max(0,\, U - I_{c,b,t})$
\end{algorithmic}
\end{algorithm}

\vspace{1.2em}
\noindent\textbf{Algorithm 4: FEFO Inventory Alert Generation}

\begin{algorithm}[H]
\caption{First-Expiry-First-Out (FEFO) Alert Generation}
\begin{algorithmic}[1]
\Require Inventory unit set $\mathcal{I}$, Alert threshold $\tau = 3$ days
\Ensure Alert queue $\mathcal{A}$ sorted by ascending remaining shelf life
\State $\mathcal{A} \leftarrow \emptyset$
\For{each unit $u \in \mathcal{I}$}
    \If{$u.\text{dispensed} = \textit{true}$} \textbf{continue} \EndIf
    \State $\text{remaining} \leftarrow u.\text{expiryDate} - \text{today}$
    \If{$\text{remaining} \leq \tau$}
        \State $\mathcal{A} \leftarrow \mathcal{A} \cup \{(u,\, \text{remaining})\}$
    \EndIf
\EndFor
\State \Return $\mathcal{A}$ sorted by ascending remaining
\end{algorithmic}
\end{algorithm}

\newpage
\begin{center}
    \textbf{\large CHAPTER 4}\\[0.4em]
    \textbf{\large IMPLEMENTATION}
\end{center}
\vspace{1.5em}

\noindent This chapter documents the complete implementation of the K9Hope platform, covering requirement analysis, input and output data specifications, the step-by-step implementation process for each functional module, hardware and software resource requirements, and the detailed design decisions made during development. All implementation work was carried out as a final-year undergraduate research project in the Department of Computer Science and Engineering, Rajalakshmi Institute of Technology, Chennai, in academic collaboration with Madras Veterinary College (MVC), Vepery. The system has not been clinically deployed; all performance metrics reported in this chapter are derived from retrospective analysis and simulation.

\vspace{2em}
\noindent\textbf{4.1\quad REQUIREMENT ANALYSIS AND SPECIFICATION}
\vspace{0.8em}

\noindent The requirement analysis for K9Hope was conducted through two complementary channels: a structured review of the DAHD July 2025 Standard Operating Procedures for Animal Blood Transfusion and Blood Banks in India, and an examination of the coordination failures documented in the IEEE platform architecture paper derived from this project. The analysis identified four categories of functional requirements: donor eligibility enforcement, blood request triage and matching, perishable inventory management, and probabilistic demand forecasting. Non-functional requirements were identified as: real-time data synchronisation across geographically distributed clinic nodes, role-based access control separating donor owners from veterinary clinics and coordinators, secure storage of medical documents, and a human-in-the-loop architecture that prevents autonomous clinical action.

\vspace{1.2em}
\noindent\textbf{4.1.1\quad Software Requirements Specification}
\vspace{0.5em}

\noindent The Software Requirements Specification (SRS) for K9Hope is organised into functional and non-functional categories.

\noindent\textbf{Functional Requirements:}

\begin{enumerate}[leftmargin=1.8em, itemsep=0.4em]
    \item \textbf{FR-01 — Donor Registration:} The system shall allow authenticated Donor Owner users to register a canine donor profile with the following mandatory fields: dog name, breed, date of birth, body weight (kg), sex, DEA blood type, most recent PCV value (\%), pathogen screening results (Babesia, Ehrlichia, Leishmania, Heartworm), vaccination history, and GPS coordinates of the owner's location. Owner consent shall be captured as a digital confirmation with a server-side timestamp.
    \item \textbf{FR-02 — MCEF Eligibility Validation:} The system shall evaluate each submitted donor profile against DAHD 2025 eligibility constraints — weight $\geq$ 25 kg, age 1--8 years, PCV $\geq$ 35\%, pathogen screens negative, and donation cooldown $\geq$ 30 days — before writing the record to the database. Profiles failing any constraint shall be rejected with a field-specific error message.
    \item \textbf{FR-03 — Blood Request Submission:} The system shall allow authenticated Veterinary Clinic users to submit blood transfusion requests specifying: patient DEA blood type required, component type (whole blood / packed RBCs / FFP / platelets / cryoprecipitate), clinical indication (free text), and an uploaded veterinary recommendation letter (PDF or JPEG/PNG, maximum 10 MB).
    \item \textbf{FR-04 — NLP Triage Classification:} The system shall automatically process uploaded recommendation letters using OCR and keyword matching, assign an urgency tier (Critical / High / Standard) based on a predefined lexicon, and surface the classified request with highlighted urgency terms to the Coordinator dashboard.
    \item \textbf{FR-05 — Geospatial Donor Matching:} Upon Coordinator approval of a blood request, the system shall query the donor registry for compatible, eligible donors within a configurable proximity radius (default 10 km) and return a shortlist sorted by ascending Haversine distance.
    \item \textbf{FR-06 — Secure Coordinator Communication:} The system shall open a request-scoped, authenticated messaging channel between Coordinator and selected Donor Owner accounts upon Coordinator-initiated outreach. No automated messages shall be sent without Coordinator action.
    \item \textbf{FR-07 — FEFO Inventory Tracking:} The system shall maintain unit-level blood component records, compute remaining shelf life daily, and generate FEFO alerts for units with remaining shelf life $\leq$ 3 days, sorted by ascending remaining life.
    \item \textbf{FR-08 — ZICP Demand Forecasting:} The system shall compute weekly prediction intervals for daily transfusion demand at each clinic using the ZICP method, and display forecast charts with upper and lower bounds on the Coordinator dashboard.
    \item \textbf{FR-09 — Audit Logging:} All Coordinator actions (request approvals, outreach initiations, inventory dispensals) shall be logged with timestamps and user IDs to a Firestore audit collection.
    \item \textbf{FR-10 — Role-Based Dashboard Routing:} On authentication, the system shall route users to a role-specific dashboard (Donor Owner / Veterinary Clinic / Coordinator) based on the role claim in the Firebase Authentication JWT token.
\end{enumerate}

\noindent\textbf{Non-Functional Requirements:}

\begin{enumerate}[leftmargin=1.8em, itemsep=0.4em]
    \item \textbf{NFR-01 — Real-Time Synchronisation:} Firestore listeners on the Coordinator dashboard shall reflect new blood requests and donor status changes within 2 seconds of database write.
    \item \textbf{NFR-02 — Security:} All API routes shall validate Firebase Authentication JWT tokens server-side. Medical documents shall be stored in Uploadcare with time-limited signed URLs. Firestore Security Rules shall enforce that each user can only read/write records associated with their own \texttt{userID} or role.
    \item \textbf{NFR-03 — Availability:} The platform shall be deployed on Vercel with edge caching, targeting $\geq$ 99.5\% uptime for the web application layer.
    \item \textbf{NFR-04 — Usability:} The interface shall be mobile-first and operable on screens of width $\geq$ 375 px without horizontal scrolling. Critical workflows (request submission, donor registration) shall be completable in $\leq$ 5 form steps.
    \item \textbf{NFR-05 — Regulatory Compliance:} The platform shall enforce DAHD 2025 SOPs programmatically and shall not facilitate commercial transactions for blood. Donor consent revocation shall be possible at any time through the Donor Owner dashboard.
\end{enumerate}

\vspace{1.2em}
\noindent\textbf{4.1.2\quad Input Data}
\vspace{0.5em}

\noindent Table 4.1 summarises the primary input data categories processed by the K9Hope platform.

\begin{center}
\textbf{Table 4.1: Input Data Specification}
\vspace{0.4em}
\begin{tabular}{p{3.2cm} p{4.5cm} p{2.2cm} p{3.5cm}}
\toprule
\textbf{Input Category} & \textbf{Fields} & \textbf{Format} & \textbf{Consuming Module} \\
\midrule
Donor Registration & Name, breed, DOB, weight, sex, DEA type, PCV, pathogen screens, GPS coords, consent & JSON form payload & MCEF, Donor Registry \\
\addlinespace
Recommendation Letter & Scanned/typed veterinary letter & PDF / JPEG / PNG & NLP Triage Module \\
\addlinespace
Blood Request & Patient DEA type, component type, clinical indication, clinic GPS coords & JSON form payload & GMA, Request Queue \\
\addlinespace
Donation Record & Donor ID, component type, volume collected, collection date & JSON (Coordinator input) & Inventory Manager, Donor Registry \\
\addlinespace
Historical Demand & Daily transfusion count per clinic per blood type, date & Firestore collection & ZICP Forecasting Engine \\
\addlinespace
Inventory Unit & Component type, DEA type, collection date, expiry date, leukoreduction flag & JSON (Coordinator input) & FEFO Alert Module \\
\bottomrule
\end{tabular}
\end{center}

\vspace{1.2em}
\noindent\textbf{4.1.3\quad Output Data}
\vspace{0.5em}

\noindent Table 4.2 summarises the primary output data produced by the platform and their consumers.

\begin{center}
\textbf{Table 4.2: Output Data Specification}
\vspace{0.4em}
\begin{tabular}{p{3.5cm} p{4.2cm} p{2.5cm} p{3cm}}
\toprule
\textbf{Output} & \textbf{Content} & \textbf{Consumer} & \textbf{Producing Module} \\
\midrule
MCEF Decision & Pass / Fail with reason code & Donor Owner, Firestore & MCEF \\
\addlinespace
NLP Urgency Classification & Tier (Critical/High/Standard), extracted keywords & Coordinator Dashboard & NLP Triage \\
\addlinespace
GMA Donor Shortlist & Ranked list of donor profiles with distances (km) & Coordinator Dashboard & GMA \\
\addlinespace
ZICP Prediction Interval & $[L, U]$ per clinic per blood type per day (7-day horizon) & Coordinator Dashboard & ZICP Engine \\
\addlinespace
FEFO Alerts & Unit ID, component type, remaining shelf life (days) & Coordinator Dashboard & FEFO Module \\
\addlinespace
Audit Log Entry & Action type, user ID, timestamp, affected record ID & Firestore audit collection & All modules \\
\addlinespace
Message Thread & Timestamped messages between Coordinator and Donor Owner & Both parties' dashboards & Communication Layer \\
\bottomrule
\end{tabular}
\end{center}

\vspace{2em}
\noindent\textbf{4.2\quad IMPLEMENTATION PROCESS AND EXPLANATION}
\vspace{0.8em}

\noindent The K9Hope platform was implemented in five sequential phases over the course of the final-year project: (1) Firebase project setup and authentication scaffolding, (2) donor registry and MCEF implementation, (3) blood request, NLP triage, and GMA implementation, (4) ZICP forecasting engine development, and (5) inventory management, FEFO alerts, and dashboard integration. Each phase is described below.

\noindent\textbf{Phase 1: Project Scaffolding and Authentication}
\vspace{0.5em}

\noindent The project was initialised using \texttt{create-next-app} with the TypeScript template and Next.js 14's App Router. Firebase was configured with a dedicated project for K9Hope, with Firestore, Authentication, Realtime Database, and Storage services enabled. Firebase Authentication was configured to support Email/Password and Google OAuth sign-in methods. A custom \texttt{authMiddleware.ts} was implemented in Next.js middleware to intercept all protected API routes, verify the Firebase ID token from the \texttt{Authorization: Bearer} header using the Firebase Admin SDK, and extract the user's \texttt{uid} and \texttt{role} custom claim. Role assignment is performed at account creation: the user selects a role (Donor Owner / Veterinary Clinic / Coordinator) during sign-up, and the role is written as a custom claim via a Cloud Function triggered on user creation.

\noindent Firestore Security Rules were written to enforce the following access control policy: Donor Owners may read and write only their own \texttt{DonorProfile} documents; Veterinary Clinics may write to \texttt{BloodRequest} and read their own requests; Coordinators may read all collections and write approval status fields. Medical documents are stored in Uploadcare rather than Firebase Storage to avoid egress cost at scale; a \texttt{/api/upload} route generates signed Uploadcare upload tokens with a 10-minute expiry and 10 MB file size limit.

\vspace{1.2em}
\noindent\textbf{Phase 2: Donor Registry and MCEF Implementation}
\vspace{0.5em}

\noindent The donor registration form was built as a multi-step React component using \texttt{react-hook-form} for client-side form state and \texttt{zod} for schema validation. The MCEF logic was implemented as a pure TypeScript function on the API route \texttt{/api/donors/register} (Next.js Route Handler), executing server-side before any Firestore write. The MCEF function evaluates constraints in the order defined in Algorithm 1 (Chapter 3) using short-circuit evaluation: each failing constraint immediately returns a structured error object with a \texttt{field} key and \texttt{reason} string, which is mapped to a field-level error on the registration form.

\noindent On MCEF pass, the donor profile is written to Firestore with the following document structure:

\begin{lstlisting}[language=Java, caption={Firestore DonorProfile document structure}]
{
  uid: string,            // Firebase Auth UID of owner
  dogName: string,
  breed: string,
  dob: Timestamp,         // Firebase Timestamp
  weightKg: number,
  sex: "M" | "F",
  deaType: string,        // e.g. "DEA1.1+", "DEA1.1-"
  pcvPercent: number,
  pathogenScreen: {
    babesia: "negative" | "positive",
    ehrlichia: "negative" | "positive",
    leishmania: "negative" | "positive",
    heartworm: "negative" | "positive"
  },
  consentTimestamp: Timestamp,
  eligibilityStatus: "verified" | "ineligible",
  lastDonationDate: Timestamp | null,
  cooldownUntil: Timestamp | null,
  locationGeo: GeoPoint,  // Firestore GeoPoint
  createdAt: Timestamp
}
\end{lstlisting}

\noindent The \texttt{cooldownUntil} field is updated by the \texttt{/api/donations/record} API route whenever a donation event is logged by a Coordinator, setting it to 30 days after the donation date. The GMA query reads this field dynamically at match time to exclude donors currently in cooldown.

\vspace{1.2em}
\noindent\textbf{Phase 3: Blood Request, NLP Triage, and GMA Implementation}
\vspace{0.5em}

\noindent\textit{Blood Request Submission.} The blood request form was implemented as a two-step form: Step 1 collects patient and request details; Step 2 handles document upload via the Uploadcare widget and triggers document upload to the signed URL. On form submission, the API route \texttt{/api/requests/submit} writes a \texttt{BloodRequest} document to Firestore and asynchronously invokes the NLP triage pipeline on the uploaded document URL.

\noindent\textit{NLP Triage Module.} The NLP triage pipeline is implemented in TypeScript on the server side. On invocation, the module fetches the uploaded document from its Uploadcare URL, extracts text using the \texttt{pdf-parse} library (for PDF inputs) or \texttt{tesseract.js} (for image inputs). The extracted plain text is then scanned using a case-insensitive keyword matching function against the three-tier urgency lexicon defined in Chapter 3. Each matched keyword is recorded with its tier and the surrounding sentence context (±30 characters) for display on the Coordinator dashboard. The highest matched tier determines the urgency classification. The classification result and matched keywords array are written back to the \texttt{BloodRequest} document in Firestore, triggering a real-time update on the Coordinator dashboard via a Firestore \texttt{onSnapshot} listener.

\noindent The urgency lexicon is defined as a TypeScript constant and is extensible without code changes to the triage logic:

\begin{lstlisting}[language=Java, caption={NLP urgency lexicon (excerpt)}]
const URGENCY_LEXICON = {
  CRITICAL: [
    "trauma", "hemoperitoneum", "acute haemorrhage",
    "haemodynamic instability", "collapse", "shock",
    "pcv < 15", "hb < 5", "gcs"
  ],
  HIGH: [
    "severe anaemia", "imha", "immune-mediated haemolytic",
    "surgical blood loss", "transfusion-dependent",
    "pcv 15", "pcv 16", "pcv 17", "pcv 18", "pcv 19",
    "pcv 20", "pcv 21", "pcv 22", "pcv 23", "pcv 24", "pcv 25"
  ],
  STANDARD: [
    "elective", "pre-operative", "chronic anaemia", "monitoring"
  ]
};
\end{lstlisting}

\noindent\textit{Geospatial Matching Algorithm.} The GMA is implemented as a Firestore query followed by a server-side distance filter. Since Firestore does not natively support geo-radius queries with combined field filters, the implementation uses a bounding-box pre-filter to reduce the candidate set before applying the exact Haversine computation. The bounding-box is computed from the clinic's GPS coordinates and the configured radius using the approximation: $\Delta\text{lat} = r / 111$ degrees, $\Delta\text{lon} = r / (111 \cdot \cos\phi_c)$ degrees. The Firestore query filters donors by \texttt{locationGeo} within the bounding box and by \texttt{deaType} and \texttt{eligibilityStatus == "verified"}. The resulting candidate set (typically 5--30 donors) is then passed through the exact Haversine computation in TypeScript. Donors with \texttt{cooldownUntil > today} are excluded at this stage. The final shortlist is sorted by distance and returned to the Coordinator dashboard.

\vspace{1.2em}
\noindent\textbf{Phase 4: ZICP Forecasting Engine Implementation}
\vspace{0.5em}

\noindent The ZICP engine is implemented in Python 3.10 as a scheduled Firebase Cloud Function that runs daily at 23:00 IST via a Cloud Scheduler trigger. The function reads the previous 90 days of \texttt{DemandRecord} documents from Firestore using the Firebase Admin SDK for Python, assembles a clinic-day-bloodtype feature matrix, and executes the ZICP training, calibration, and inference pipeline described in Algorithm 3 (Chapter 3).

\noindent The data split follows a chronological protocol consistent with the experimental setup reported in the IEEE ZICP paper: 50\% training, 25\% calibration, 25\% held-out test (for evaluation); at inference, the full 75\% non-test window is used for training and calibration combined, with the most recent 25\% serving as the calibration set. Feature engineering extracts: day-of-week (one-hot, 7 dimensions), month (one-hot, 12 dimensions), clinic ID (label encoded), blood type (label encoded), and a 7-day rolling mean of demand for the same clinic-blood-type combination.

\noindent The occurrence model (\texttt{LGBMClassifier}) and size model (\texttt{LGBMRegressor}) are trained with the following hyperparameters, selected by 3-fold time-series cross-validation on the training set:

\begin{center}
\textbf{Table 4.3: ZICP LightGBM Hyperparameters}
\vspace{0.4em}
\begin{tabular}{lll}
\toprule
\textbf{Parameter} & \textbf{Occurrence Model} & \textbf{Size Model} \\
\midrule
\texttt{n\_estimators} & 200 & 200 \\
\texttt{learning\_rate} & 0.05 & 0.05 \\
\texttt{max\_depth} & 4 & 4 \\
\texttt{num\_leaves} & 15 & 15 \\
\texttt{min\_child\_samples} & 10 & 5 \\
\texttt{subsample} & 0.8 & 0.8 \\
\texttt{colsample\_bytree} & 0.8 & 0.8 \\
\texttt{objective} & \texttt{binary} & \texttt{regression} \\
\texttt{random\_state} & 42 & 42 \\
\bottomrule
\end{tabular}
\end{center}

\noindent The conformal calibration procedure computes non-conformity scores on the calibration set and extracts the $\lceil(n_\text{cal}+1)(1-\alpha)\rceil/n_\text{cal}$-th order statistic as the calibrated threshold $q$. At inference, the prediction interval upper bound $U$ is found by binary search over $d \in \{0, 1, 2, \ldots, 20\}$, finding the maximum $d$ such that $S(X_\text{new}, d) \leq q$. The lower bound is always $L = 0$, since demand is non-negative and the ZICP interval includes zero by construction. The resulting 7-day-ahead prediction intervals $\{[L_{c,b,t+k}]_{k=1}^{7}\}$ are written to a \texttt{Forecasts} Firestore collection keyed by clinic ID, blood type, and forecast date, and visualised as shaded band charts on the Coordinator dashboard using the Recharts library.

\vspace{1.2em}
\noindent\textbf{Phase 5: Inventory Management, FEFO Alerts, and Dashboard Integration}
\vspace{0.5em}

\noindent\textit{Inventory Unit Management.} Coordinators log incoming blood units through a form on the inventory panel. Each submitted unit creates an \texttt{InventoryUnit} Firestore document with fields as defined in Section 4.1.2. Expiry dates are auto-populated based on component type according to the shelf life rules defined in Table 4.4.

\begin{center}
\textbf{Table 4.4: Blood Component Shelf Life Rules}
\vspace{0.4em}
\begin{tabular}{lll}
\toprule
\textbf{Component Type} & \textbf{Storage Condition} & \textbf{Shelf Life} \\
\midrule
Fresh Whole Blood & Unrefrigerated & 8 hours \\
Stored Whole Blood & $1$--$6^\circ$C & 28--35 days \\
Packed Red Blood Cells & $1$--$6^\circ$C & 28--35 days \\
Fresh Frozen Plasma & $-18^\circ$C & 12 months \\
Platelets & $20$--$24^\circ$C (agitated) & 5 days \\
Cryoprecipitate & $-18^\circ$C & 12 months \\
\bottomrule
\end{tabular}
\end{center}

\noindent\textit{FEFO Alert Generation.} A daily Cloud Function queries all non-dispensed \texttt{InventoryUnit} documents and computes remaining shelf life as \texttt{expiryDate - today}. Units with remaining life $\leq 3$ days are written to a \texttt{FEFOAlerts} Firestore collection. The Coordinator dashboard subscribes to this collection via a Firestore \texttt{onSnapshot} listener and renders alerts as a colour-coded table (red for $\leq 1$ day, amber for 2--3 days).

\noindent\textit{Dashboard Integration.} The Coordinator dashboard aggregates five real-time data streams: the NLP-triaged request queue (from \texttt{BloodRequests} collection), GMA shortlist results (on-demand, triggered by Coordinator approval), ZICP forecast charts (from \texttt{Forecasts} collection), FEFO alerts (from \texttt{FEFOAlerts} collection), and the messaging panel (from Firebase Realtime Database). Each stream is managed by a dedicated React custom hook using the \texttt{useEffect} / \texttt{onSnapshot} pattern to subscribe and unsubscribe on component mount/unmount, preventing memory leaks during navigation between dashboard sections.

\vspace{2em}
\noindent\textbf{4.3\quad RESOURCE REQUIREMENTS}
\vspace{0.8em}

\noindent This section specifies the hardware and software resources used in the development and evaluation of the K9Hope platform.

\vspace{1.2em}
\noindent\textbf{4.3.1\quad Hardware}
\vspace{0.5em}

\noindent The platform was developed and the ZICP model was trained and evaluated on the following hardware configuration, consistent with the experimental setup reported in the IEEE platform architecture paper.

\begin{center}
\textbf{Table 4.5: Development Hardware Specification}
\vspace{0.4em}
\begin{tabular}{ll}
\toprule
\textbf{Component} & \textbf{Specification} \\
\midrule
Processor & Intel Core i7-12700H (12th Gen, 14 cores, up to 4.7 GHz) \\
RAM & 32 GB DDR5 \\
Storage & 512 GB NVMe SSD \\
Operating System & Windows 11 Pro / Ubuntu 22.04 LTS (WSL2) \\
Network & Standard broadband ($\geq$ 50 Mbps) for Firebase and Vercel deployments \\
\bottomrule
\end{tabular}
\end{center}

\noindent For production deployment, the platform operates entirely on cloud infrastructure. The Next.js frontend and API routes are deployed on Vercel's serverless edge network; Firebase Cloud Functions run on Google Cloud's managed Node.js 20 and Python 3.10 runtimes. No dedicated server hardware is required for production operation. The minimum client hardware for accessing the web application is a smartphone with a modern browser (Chrome 90+, Safari 14+) and a 4G or Wi-Fi connection.

\vspace{1.2em}
\noindent\textbf{4.3.2\quad Software}
\vspace{0.5em}

\noindent Table 4.6 lists the complete software stack used in the development, testing, and deployment of the K9Hope platform.

\begin{center}
\textbf{Table 4.6: Software Stack}
\vspace{0.4em}
\begin{tabular}{p{3.5cm} p{3cm} p{6cm}}
\toprule
\textbf{Component} & \textbf{Version} & \textbf{Purpose} \\
\midrule
Next.js & 14.2 & Frontend framework with App Router and server-side rendering \\
TypeScript & 5.3 & Type-safe frontend and API route development \\
Tailwind CSS & 3.4 & Utility-first responsive styling \\
shadcn/ui & 0.8 & Accessible React component library \\
Recharts & 2.10 & ZICP forecast chart visualisation \\
react-hook-form & 7.49 & Multi-step form state management \\
zod & 3.22 & Schema validation for API inputs \\
Firebase Admin SDK & 12.0 (Node.js) & Server-side Firestore, Auth, and Realtime DB access \\
Firebase Client SDK & 10.7 & Client-side real-time listeners and Auth \\
Firebase Firestore & — & NoSQL real-time database for all persistent records \\
Firebase Authentication & — & JWT-based role-differentiated access control \\
Firebase Realtime Database & — & Low-latency messaging between Coordinator and Donor Owner \\
Uploadcare & — & Secure CDN storage for medical documents \\
Vercel & — & Cloud deployment, serverless API routes, edge caching \\
Python & 3.10.12 & ZICP forecasting engine and data processing \\
LightGBM & 4.0.0 & Gradient boosting for occurrence and size models \\
scikit-learn & 1.3.0 & Preprocessing, cross-validation, SMOTE \\
pandas & 2.0.3 & Demand data manipulation and feature engineering \\
numpy & 1.25.2 & Numerical computation for ZICP calibration \\
pdf-parse & 1.1.1 & PDF text extraction for NLP triage \\
tesseract.js & 5.0 & OCR for image-format recommendation letters \\
Node.js & 20 LTS & Runtime for Firebase Cloud Functions \\
Git / GitHub & — & Version control and team collaboration \\
Visual Studio Code & — & Primary development IDE \\
\bottomrule
\end{tabular}
\end{center}

\vspace{2em}
\noindent\textbf{4.4\quad DESIGN}
\vspace{0.8em}

\noindent This section documents the key design decisions made during implementation, covering database schema design, API design, UI/UX design, and the design of the ZICP data pipeline.

\vspace{1.2em}
\noindent\textbf{4.4.1\quad Design Process and Key Design Decisions}
\vspace{0.5em}

\noindent\textit{Database Schema Design.} Firestore's document-collection model was selected over a relational database to support real-time multi-client synchronisation without polling. The schema uses shallow, denormalised documents to minimise read latency on the Coordinator dashboard: donor contact details are embedded in the \texttt{DonorProfile} document rather than referenced through a separate \texttt{User} document join, allowing the GMA shortlist query to return complete donor information in a single Firestore read. The \texttt{InventoryUnit} and \texttt{DemandRecord} collections use composite document IDs of the format \texttt{\{clinicID\}\_\{bloodType\}\_\{date\}} to enable efficient single-document reads and avoid expensive collection-group queries.

\noindent\textit{API Design.} All backend logic is encapsulated in Next.js Route Handlers under the \texttt{/api/} directory, following a RESTful pattern. The four primary API endpoints are:

\begin{itemize}[leftmargin=1.5em, itemsep=0.3em]
    \item \texttt{POST /api/donors/register} — Runs MCEF, writes to DonorProfile collection
    \item \texttt{POST /api/requests/submit} — Writes BloodRequest, triggers NLP triage async
    \item \texttt{POST /api/requests/[id]/approve} — Coordinator approval, triggers GMA
    \item \texttt{POST /api/donations/record} — Logs donation event, updates cooldownUntil
\end{itemize}

\noindent All routes validate the Firebase ID token, extract the caller's role, and reject requests where the caller's role does not match the required role for the operation (e.g., only Coordinators may call \texttt{/api/requests/[id]/approve}).

\noindent\textit{UI/UX Design.} The interface design followed a mobile-first approach using Tailwind CSS breakpoints. The Coordinator dashboard was designed as a tabbed single-page application with four panels: Request Queue, Inventory, Forecast, and Messages. Each panel uses a card-based layout with colour-coded status indicators — red for Critical urgency and near-expiry alerts, amber for High urgency and 2--3 day shelf life, green for Standard urgency and healthy inventory. The donor registration form uses a three-step wizard pattern (Basic Info $\rightarrow$ Health Details $\rightarrow$ Consent \& Submit) to reduce cognitive load for pet owners who may be unfamiliar with medical terminology.

\noindent\textit{ZICP Data Pipeline Design.} The ZICP pipeline was designed to operate correctly on very small clinic-level datasets. Individual clinics in the retrospective dataset recorded fewer than 200 transfusion events per year, with 78\% zero-demand days. To prevent data leakage, the pipeline strictly respects chronological ordering: training precedes calibration which precedes test in all experiments. SMOTE oversampling, applied in the donor return classification experiments in the IEEE platform architecture paper, was not used in ZICP training since the task is a regression on occurrence probability, not classification. The LightGBM \texttt{min\_child\_samples=10} hyperparameter was set conservatively to prevent leaf-level overfitting on the positive-demand records, which constitute only 22\% of the full dataset.

\noindent\textit{Human-in-the-Loop Design.} Every automated output in the system — eligibility decisions, urgency classifications, donor shortlists, forecast intervals, and expiry alerts — is presented to a Coordinator as a recommendation, not an action. The UI reinforces this through explicit ``Approve'' and ``Initiate Outreach'' action buttons that require Coordinator interaction before any downstream effect occurs. This design choice is directly motivated by the ethical constraints of canine blood donation documented in the IEEE papers: owner consent is revocable at any time, automated contact risks donor welfare, and the epistemic uncertainty of models trained on sparse data is too high to justify autonomous action.

\newpage
\begin{center}
    \textbf{\large CHAPTER 5}\\[0.4em]
    \textbf{\large RESULTS AND DISCUSSION}
\end{center}
\vspace{1.5em}

\noindent This chapter presents the experimental evaluation of the K9Hope platform across all primary performance metrics. Evaluation was conducted through two complementary methods: (1) retrospective analysis on 1,847 transfusion records collected from 12 veterinary clinics over a two-year period (January 2022 -- December 2023), and (2) simulation on 5,000 synthetic donor profiles processed against 500 emergency blood request scenarios. All results reported here are derived from these offline evaluations; no prospective clinical deployment has been undertaken. The results are discussed in the context of their operational significance, and the limitations of the simulation-based methodology are explicitly acknowledged.

\vspace{2em}
\noindent\textbf{5.1\quad MCEF ELIGIBILITY FILTER PERFORMANCE}
\vspace{0.8em}

\noindent The MCEF was evaluated on 5,000 synthetic donor profiles generated to include a range of valid and invalid attribute combinations across all six eligibility constraints. Ground-truth eligibility labels were assigned by applying the DAHD 2025 criteria analytically to each profile. The MCEF achieved a constraint classification accuracy of 100\% across all 5,000 profiles: every profile that violated one or more constraints was correctly rejected, and every profile satisfying all constraints was correctly admitted. No false negatives (ineligible donors passed as eligible) were observed. This result reflects the deterministic nature of the MCEF: it implements Boolean constraint checks on structured numerical inputs with no probabilistic component, so perfect accuracy is expected given well-formed inputs.

\noindent Table 5.1 shows the distribution of constraint violations across the synthetic profile set.

\begin{center}
\textbf{Table 5.1: MCEF Constraint Violation Distribution (5,000 Synthetic Profiles)}
\vspace{0.4em}
\begin{tabular}{lrr}
\toprule
\textbf{Violated Constraint} & \textbf{Profiles Rejected} & \textbf{\% of Total} \\
\midrule
Weight $<$ 25 kg & 1,142 & 22.8\% \\
Age out of range (1--8 years) & 893 & 17.9\% \\
PCV $<$ 35\% & 674 & 13.5\% \\
Pathogen screen positive & 521 & 10.4\% \\
Within 30-day cooldown & 318 & 6.4\% \\
Owner consent not recorded & 112 & 2.2\% \\
Multiple constraints violated & 487 & 9.7\% \\
\midrule
Total rejected & 3,653 & 73.1\% \\
Total admitted (verified) & 1,347 & 26.9\% \\
\bottomrule
\end{tabular}
\end{center}

\noindent The admission rate of 26.9\% reflects the restrictive biological and regulatory eligibility criteria under DAHD 2025 SOPs, consistent with the observation in the IEEE platform architecture paper that the small fraction of owned dogs eligible to donate is a primary structural constraint on canine blood service capacity. The MCEF eliminates the need for manual eligibility verification at the time of emergency, ensuring that all donors in the registry are pre-screened and immediately actionable when a blood request is received.

\vspace{2em}
\noindent\textbf{5.2\quad GEOSPATIAL MATCHING — TIME-TO-MATCH AND SUCCESS RATE}
\vspace{0.8em}

\noindent The GMA was evaluated by simulating 500 emergency blood requests against the pool of 1,347 verified synthetic donor profiles, with clinic and donor locations distributed across a simulated Chennai metropolitan region. Each request was assigned a required DEA blood type drawn from the empirical blood type distribution observed in the retrospective clinical dataset. The 10 km default radius was used for all simulations. Table 5.2 compares K9Hope's simulated performance against the manual baseline figures documented in the retrospective clinical data.

\begin{center}
\textbf{Table 5.2: Donor Matching Performance — K9Hope vs. Manual Baseline}
\vspace{0.4em}
\begin{tabular}{lll}
\toprule
\textbf{Metric} & \textbf{Manual Baseline} & \textbf{K9Hope (Simulated)} \\
\midrule
Median time-to-match & 240 minutes & 12 minutes \\
Donor identification success rate (4-hour window) & 60\% & 98\% \\
Median donor travel distance & $>$ 20 km & $<$ 10 km \\
Donors rejected on arrival (ineligible) & High (unquantified) & 0\% (pre-screened) \\
\bottomrule
\end{tabular}
\end{center}

\noindent The 12-minute median time-to-match reflects the GMA's ability to execute a proximity-filtered, eligibility-checked donor query in under one second, followed by Coordinator review and outreach initiation. The dominant component of the 12-minute figure is the Coordinator's review and communication time, not computational latency. The 98\% success rate within the 4-hour critical response window represents a substantial improvement over the 60\% baseline, and is primarily attributable to the pre-screening benefit of the MCEF: in the manual system, a significant proportion of the 240-minute median search time was consumed by contacting donors who were subsequently found to be ineligible upon arrival. By eliminating this waste, K9Hope compresses the effective search space to verified, immediately actionable donors.

\noindent It is important to note that the 98\% success rate is a simulation result and depends on the density of verified donors within the 10 km radius of each requesting clinic. In sparsely populated regions or for rare DEA blood types, the success rate will be lower. The configurable radius parameter allows coordinators to expand the search boundary in such cases.

\vspace{2em}
\noindent\textbf{5.3\quad NLP TRIAGE CLASSIFICATION PERFORMANCE}
\vspace{0.8em}

\noindent The NLP triage module was evaluated on a labelled set of 200 veterinary recommendation letters extracted from the retrospective clinical dataset, manually annotated by a veterinary professional at Madras Veterinary College with ground-truth urgency tier labels (Critical / High / Standard). The module's keyword-based classification was compared against these labels. Table 5.3 presents the classification performance.

\begin{center}
\textbf{Table 5.3: NLP Triage Classification Performance (200 Annotated Letters)}
\vspace{0.4em}
\begin{tabular}{lrrrr}
\toprule
\textbf{Urgency Tier} & \textbf{Precision} & \textbf{Recall} & \textbf{F1-Score} & \textbf{Support} \\
\midrule
Critical & 0.91 & 0.88 & 0.89 & 67 \\
High & 0.84 & 0.86 & 0.85 & 79 \\
Standard & 0.92 & 0.94 & 0.93 & 54 \\
\midrule
Macro Average & 0.89 & 0.89 & 0.89 & 200 \\
\bottomrule
\end{tabular}
\end{center}

\noindent The macro-average F1-score of 0.89 is satisfactory for a keyword-based approach given the heterogeneity of clinical writing styles in the letter corpus. The most common error mode was Critical letters being classified as High (8 cases), typically where the document described critical haemodynamic parameters using non-standard abbreviations or Tamil-language terms not captured in the English keyword lexicon. This limitation is acknowledged: the NLP module's performance degrades on handwritten or non-standardised documents, and OCR quality significantly affects accuracy for scanned letters. The module's purpose is triage support, not diagnosis; all classifications are reviewed by a Coordinator before any action is taken, which mitigates the impact of classification errors.

\vspace{2em}
\noindent\textbf{5.4\quad ZICP DEMAND FORECASTING PERFORMANCE}
\vspace{0.8em}

\noindent The ZICP forecasting engine was evaluated on the 12-clinic retrospective dataset using a chronological 50/25/25 train/calibration/test split, following the experimental protocol described in the IEEE ZICP paper. Results were averaged over 5 random seeds. Table 5.4 presents the full performance comparison across methods at significance level $\alpha = 0.05$.

\begin{center}
\textbf{Table 5.4: Demand Forecasting Performance — Method Comparison ($\alpha = 0.05$)}
\vspace{0.4em}
\begin{tabular}{lrrrr}
\toprule
\textbf{Method} & \textbf{Coverage} & \textbf{Interval Width} & \textbf{Interval Score} & \textbf{Inventory Cost (INR)} \\
\midrule
Naive (historical quantiles) & 0.89 & 5.2 & 8.3 & 1,240 \\
Bootstrap (quantile regression) & 0.91 & 4.8 & 7.1 & 1,180 \\
Bayesian NN (MC dropout) & 0.87 & 6.1 & 9.2 & 1,350 \\
ARIMA$+2\sigma$ & 0.84 & 5.8 & 8.9 & 1,290 \\
\textbf{ZICP (proposed)} & \textbf{0.92} & \textbf{3.4} & \textbf{5.2} & \textbf{970} \\
\bottomrule
\end{tabular}
\end{center}

\noindent ZICP achieves 92\% empirical coverage at the 95\% target, the closest of all evaluated methods to the nominal level, while producing prediction intervals that are 35\% tighter than the Bootstrap baseline (width 3.4 vs. 4.8 units) and 44\% tighter than the Bayesian NN. The interval score of 5.2 represents a 27\% improvement over the next-best method (Bootstrap, 7.1), confirming that ZICP's advantage is in calibration efficiency rather than simply wider intervals. The inventory ordering cost of INR 970 per clinic-week represents a 22\% reduction relative to the Bootstrap baseline (INR 1,180) and a 28\% reduction relative to the ARIMA-based safety-stock heuristic (INR 1,290).

\noindent Figure 5.1 illustrates the empirical coverage across three significance levels ($\alpha \in \{0.05, 0.10, 0.20\}$). ZICP consistently tracks closest to the target coverage line, while all other methods show systematic under-coverage at $\alpha = 0.05$ and over-coverage at $\alpha = 0.20$.

\begin{figure}[H]
    \centering
    \includegraphics[width=\linewidth]{fig_5_1_zicp_coverage}
    \caption{\textbf{Fig. 5.1 ZICP empirical coverage vs. target across significance levels (upper) and average interval width (lower). ZICP achieves the tightest intervals closest to the coverage target.}}
\end{figure}

\noindent\textbf{Ablation Study.} Table 5.5 presents the ablation results comparing ZICP-full (both occurrence and size models) against ZICP-occ (occurrence model only, size model replaced by historical mean) and ZICP-size (size model only, occurrence model replaced by constant $\hat{\pi} = 0.22$). The full model outperforms both ablations on all metrics, confirming that jointly modelling both zero-inflation components is necessary for accurate calibrated intervals on this dataset.

\begin{center}
\textbf{Table 5.5: ZICP Ablation Study}
\vspace{0.4em}
\begin{tabular}{lrrrr}
\toprule
\textbf{Variant} & \textbf{Coverage} & \textbf{Width} & \textbf{Int.\ Score} & \textbf{Cost (INR)} \\
\midrule
ZICP-occ only & 0.88 & 4.1 & 6.8 & 1,090 \\
ZICP-size only & 0.85 & 3.9 & 7.2 & 1,140 \\
\textbf{ZICP-full} & \textbf{0.92} & \textbf{3.4} & \textbf{5.2} & \textbf{970} \\
\bottomrule
\end{tabular}
\end{center}

\noindent\textbf{Discussion.} The 92\% empirical coverage is 3 percentage points below the 95\% target, attributable to finite calibration sample size variance. With $n_\text{cal} \approx 450$ clinic-day records per clinic, the standard error of empirical coverage is approximately $\sqrt{0.92 \times 0.08 / 450} \approx 0.013$, so the observed shortfall is within expected sampling variation. The exchangeability assumption underlying the coverage guarantee may be weakly violated during seasonal vector-borne disease peaks, when demand patterns shift abruptly — a limitation acknowledged in the IEEE ZICP paper. The ARIMA-family models, evaluated in the earlier IEEE platform architecture paper, detected weekly seasonality (elevated Monday--Tuesday demand, coefficient of variation $>$ 100\% on low-volume days) but systematically failed to capture emergency-driven spikes, producing MAE of 0.71 and MSE of 0.85. ZICP supersedes these baselines by providing calibrated intervals rather than point estimates, directly addressing the inventory decision trade-off between shortage risk and wastage that point forecasts cannot resolve.

\vspace{2em}
\noindent\textbf{5.5\quad INVENTORY MANAGEMENT AND WASTAGE REDUCTION}
\vspace{0.8em}

\noindent The FEFO inventory management module was evaluated by simulating 90 days of inventory operations across 12 clinics using the retrospective demand data. The simulation compared three inventory strategies: (1) static safety-stock heuristic (keep 3 units of each component type), (2) ARIMA-based point forecast ordering, and (3) ZICP upper-bound ordering. Table 5.6 presents the results.

\begin{center}
\textbf{Table 5.6: Inventory Strategy Comparison (90-Day Simulation, 12 Clinics)}
\vspace{0.4em}
\begin{tabular}{lrrr}
\toprule
\textbf{Strategy} & \textbf{Wastage Rate} & \textbf{Shortage Events} & \textbf{Weekly Cost (INR)} \\
\midrule
Static safety-stock & 18--25\% & 47 & 1,290 \\
ARIMA point forecast & 14\% & 38 & 1,180 \\
\textbf{ZICP upper-bound ordering} & \textbf{8\%} & \textbf{22} & \textbf{970} \\
\bottomrule
\end{tabular}
\end{center}

\noindent ZICP-based ordering reduces blood wastage from the 18--25\% range observed with static heuristics to 8\%, a reduction of 67\% at the lower bound, consistent with the operational significance target identified in the problem statement. The 22 shortage events over 90 days across 12 clinics (approximately 0.02 per clinic-day) remain non-zero because ZICP intervals do not guarantee zero shortages — they guarantee coverage at the specified level, meaning approximately 8\% of days will see demand exceeding the upper bound at $\alpha = 0.05$. Clinics can reduce shortage events further by using a lower significance level (e.g., $\alpha = 0.01$), at the cost of wider intervals and higher ordering quantities. The FEFO alert system, by ensuring near-expiry units are dispensed first, contributes directly to the wastage reduction by preventing units from expiring while newer units are used preferentially.

\vspace{2em}
\noindent\textbf{5.6\quad DONOR RETURN CLASSIFICATION}
\vspace{0.8em}

\noindent Donor return classification — predicting whether a registered donor will donate again within 18 months — was evaluated on 176 donors from the retrospective dataset using 5-fold stratified cross-validation with SMOTE oversampling applied only within training folds to prevent data leakage. Six classifiers were evaluated. Table 5.7 presents the results, reproduced from the IEEE platform architecture paper.

\begin{center}
\textbf{Table 5.7: Donor Return Classification Performance (Cross-Validated)}
\vspace{0.4em}
\begin{tabular}{lrrrrl}
\toprule
\textbf{Classifier} & \textbf{Accuracy} & \textbf{Precision} & \textbf{Recall} & \textbf{F1} & \textbf{$p$-value} \\
\midrule
ANN (1 hidden layer) & 78.4\% & 74.2\% & 81.6\% & 77.7\% & 0.042$^*$ \\
Random Forest & 76.5\% & 73.8\% & 68.3\% & 70.9\% & 0.041$^*$ \\
Logistic Regression & 72.1\% & 68.0\% & 70.4\% & 69.1\% & 0.089 \\
SVM & 70.3\% & 69.1\% & 67.0\% & 68.0\% & 0.156 \\
Decision Tree & 69.5\% & 63.5\% & 66.2\% & 64.8\% & 0.203 \\
Naive Bayes (baseline) & 68.8\% & 64.0\% & 66.8\% & 65.3\% & — \\
\bottomrule
\multicolumn{6}{l}{$^*$ $p < 0.05$ vs. Naive Bayes baseline (Wilcoxon Signed-Rank Test)}
\end{tabular}
\end{center}

\noindent The ANN achieved the highest accuracy (78.4\%, F1 77.7\%), with Random Forest the next best (76.5\%, F1 70.9\%). Both results are statistically significant at $p < 0.05$ relative to the Naive Bayes baseline. However, as documented transparently in the IEEE platform architecture paper, the 1.9 percentage point accuracy gap between ANN and Random Forest has no operational significance given the small sample size ($n = 176$), and the ANN almost certainly overfits: the persistent training-validation loss gap (0.32 vs. 0.58 at early stopping epoch 487) is consistent with memorisation of individual records in a single-hidden-layer network trained on fewer than 200 samples. Random Forest is the operationally recommended classifier for this task: it provides more conservative, balanced precision-recall tradeoff (specificity 82.4\%, sensitivity 68.3\%), is more interpretable, and is less susceptible to overfitting. The strongest predictors from Random Forest feature importance were months since last donation, total prior donations, and body weight. The weak association of DEA blood type with return behaviour is likely confounded by clinic calling practices: universal donors (DEA 1.1$-$) may be called more frequently, artificially inflating their apparent return rate.

\noindent The critical limitation of these results is the absence of owner-level variables — distance from clinic, socioeconomic status, communication channel preference, and emotional response to first donation — which are likely more predictive than the biological features available in administrative records. These results should not guide operational outreach prioritisation without larger validation datasets that include owner-level characteristics.

\vspace{2em}
\noindent\textbf{5.7\quad OVERALL SYSTEM EVALUATION SUMMARY}
\vspace{0.8em}

\noindent Table 5.8 consolidates the primary evaluation metrics across all platform modules.

\begin{center}
\textbf{Table 5.8: K9Hope Platform Evaluation Summary}
\vspace{0.4em}
\begin{tabular}{p{4cm} p{3.5cm} p{5.5cm}}
\toprule
\textbf{Module} & \textbf{Key Metric} & \textbf{Result} \\
\midrule
MCEF & Eligibility filter accuracy & 100\% (deterministic, structured inputs) \\
GMA & Median time-to-match & 12 min (vs. 240 min manual baseline) \\
GMA & Success rate (4-hr window) & 98\% (vs. 60\% manual baseline) \\
NLP Triage & Macro F1-score & 0.89 (200 annotated letters) \\
ZICP & Empirical coverage & 92\% (target 95\%, $\alpha = 0.05$) \\
ZICP & Interval width vs. Bootstrap & 35\% tighter (3.4 vs. 4.8 units) \\
ZICP & Inventory cost reduction & 22\% (INR 970 vs. INR 1,180 baseline) \\
FEFO + ZICP & Blood wastage reduction & 67\% lower bound (8\% vs. 18--25\%) \\
Donor Return & Best classifier F1 & 77.7\% (ANN, simulation-only) \\
\bottomrule
\end{tabular}
\end{center}

\noindent The principal finding of this evaluation is consistent with the conclusion of both IEEE papers derived from this project: the dominant barrier to canine blood service coordination in India is information fragmentation, not algorithmic complexity. The largest performance gains — the reduction in time-to-match from 240 minutes to 12 minutes and the increase in success rate from 60\% to 98\% — are attributable to the structural integration of a pre-screened donor registry with geospatial search, not to sophisticated machine learning. The ZICP method provides a genuine methodological contribution over ARIMA-family baselines by delivering calibrated prediction intervals rather than point estimates, but its operational value is also fundamentally enabled by the existence of a centralised demand record database, which currently does not exist in Indian veterinary practice. The platform's most tractable near-term impact is therefore as an information coordination infrastructure, with the analytical modules providing incremental benefit as data accumulates.

\newpage
\begin{center}
    \textbf{\large CHAPTER 6}\\[0.4em]
    \textbf{\large CONCLUSION}
\end{center}
\vspace{1.5em}

\noindent This project presented K9Hope, India's first digital decision support platform for canine blood donation coordination in sparse veterinary networks. The platform was designed, implemented, and evaluated in response to a documented operational failure: when a dog requires an emergency blood transfusion in India, there is currently no structured system to connect the patient with a compatible, eligible, and nearby donor. The manual search process takes a median of 240 minutes and succeeds within the critical response window only 60\% of the time. K9Hope directly addresses this gap through a layered, human-in-the-loop architecture integrating four algorithmic modules — the Multi-Constraint Eligibility Filter (MCEF), the Haversine-based Geospatial Matching Algorithm (GMA), the NLP document triage pipeline, and the Zero-Inflated Conformal Prediction (ZICP) demand forecasting engine — within a unified web application built on Next.js, Firebase, and Vercel.

\vspace{1.5em}
\noindent\textbf{Principal Findings.} The evaluation results establish three principal findings. First, information fragmentation is the dominant barrier to canine blood service coordination, not algorithmic complexity. The most operationally significant improvements — the reduction in time-to-match from 240 to 12 minutes and the increase in donor identification success rate from 60\% to 98\% within the 4-hour critical window — arise from structural data integration: a centralised, pre-screened donor registry searchable by blood type and proximity replaces phone calls to unverified personal contacts. This finding replicates and extends the architectural conclusion of the earlier IEEE platform architecture paper, which identified information fragmentation as the most tractable barrier and recommended a multi-clinic coordination platform as the primary intervention.

\noindent Second, the ZICP method provides a meaningful methodological advance over existing forecasting approaches for sparse veterinary demand data. By separately modelling demand occurrence and magnitude and calibrating prediction intervals using held-out data without distributional assumptions, ZICP achieves 92\% empirical coverage with intervals 35\% tighter than Bootstrap alternatives, enabling a 22\% reduction in inventory ordering costs and a reduction in blood wastage from 18--25\% (static heuristics) to 8\%. This result is practically significant because it directly supports the inventory trade-off between shortage risk and component wastage that point-estimate methods such as ARIMA cannot address. ARIMA, evaluated in the earlier paper, captures weekly seasonality but systematically fails to model the zero-inflated emergency spikes that dominate veterinary transfusion demand.

\noindent Third, donor return classification from administrative records alone produces statistically detectable but operationally limited signal. The best classifier (ANN, F1 77.7\%) is likely overfitted on the 176-donor retrospective sample, and the absence of owner-level variables — distance, socioeconomic status, communication preferences — limits the practical utility of these predictions. Random Forest (F1 70.9\%) is the recommended operational model for its balance of interpretability and conservatism, but meaningful deployment requires larger datasets that include owner-level characteristics.

\vspace{1.5em}
\noindent\textbf{Contributions.} The specific contributions of this project are as follows:

\begin{enumerate}[leftmargin=1.8em, itemsep=0.5em]
    \item The design and implementation of a complete, cloud-deployed web application that programmatically enforces DAHD 2025 donor eligibility criteria, replacing manual paper-based verification at the time of emergency.
    \item The implementation and evaluation of the Haversine-based GMA for proximity-ranked donor shortlist generation, reducing median time-to-match by 95\% relative to the manual baseline in simulation.
    \item The design of an NLP triage pipeline for veterinary recommendation letters achieving macro F1 of 0.89 on a heterogeneous annotated corpus, providing urgency-ranked request queuing for coordinators managing concurrent emergencies.
    \item The proposal, implementation, and empirical validation of the Zero-Inflated Conformal Prediction (ZICP) method for blood demand forecasting in sparse veterinary networks, providing the first distribution-free, finite-sample coverage guarantee specifically designed for zero-inflated veterinary demand data.
    \item The formalisation of DAHD 2025 animal welfare legislation within a software architecture, ensuring that consent revocability, donation cooldown enforcement, and pathogen screening compliance are structurally enforced rather than procedurally reliant on clinical staff memory.
    \item The transparent documentation of methodological limitations, including the simulation-only evaluation context, the exchangeability assumption required for ZICP coverage guarantees, the overfitting risk in donor return classification on small samples, and the absence of PMS API integration required for real-world data ingestion.
\end{enumerate}

\vspace{1.5em}
\noindent\textbf{Limitations.} The platform has not been clinically deployed. All performance metrics are derived from retrospective data analysis and simulation on synthetic profiles, and prospective validation in an operational veterinary setting remains necessary before the system can be adopted in practice. The retrospective dataset — 1,847 transfusion events from 12 clinics — is a convenience sample without systematic sampling or power calculation. The NLP triage module underperforms on handwritten or non-English-language letters. The ZICP exchangeability assumption may be violated during epidemic-driven demand shifts. Data interoperability with proprietary veterinary practice management systems (PMS) is the single most significant practical barrier to deployment: without API integration or automated ETL pipelines, manual data entry would severely limit adoption at scale.

\vspace{1.5em}
\noindent\textbf{Future Work.} The most impactful directions for future development are:

\begin{enumerate}[leftmargin=1.8em, itemsep=0.4em]
    \item \textbf{Prospective clinical deployment} in partnership with Madras Veterinary College and a minimum of three additional Chennai-area veterinary clinics, with institutional review and veterinarian usability evaluation, to validate simulation-derived metrics against real operational outcomes.
    \item \textbf{PMS API integration} with major Indian veterinary practice management platforms (e.g., VetPort, ezyVet) to enable automated ingestion of transfusion records and elimination of manual data entry, which is the primary practical barrier to adoption identified in the IEEE platform architecture paper.
    \item \textbf{Multi-step ZICP forecasting} extending the current one-week-ahead prediction to longer horizons (2--4 weeks), incorporating spatial demand correlation between geographically proximate clinics, and developing adaptive conformal methods that relax the exchangeability assumption for non-stationary demand patterns.
    \item \textbf{Multilingual NLP triage} extending the urgency lexicon to Tamil and other Indian languages, and replacing keyword matching with a fine-tuned small language model (e.g., IndicBERT) trained on a larger annotated corpus of veterinary letters, to improve classification robustness on non-standardised clinical text.
    \item \textbf{Owner-level donor retention modelling} incorporating distance from clinic, communication channel preference, and longitudinal donation history into the return classification model, requiring prospective data collection during clinical deployment.
    \item \textbf{Nationwide registry expansion} beyond the Chennai metropolitan region, potentially in partnership with the Department of Animal Husbandry and Dairying, to realise the network-effect value of a centralised national canine donor registry.
\end{enumerate}

\vspace{1.5em}
\noindent The central lesson of this project is that the gap between a dog in a life-threatening emergency and the blood that could save it is not a gap of algorithmic sophistication — it is a gap of structured information. K9Hope demonstrates that fundamental computer science tools — rule-based constraint filters, Haversine distance computation, keyword matching, and held-out calibration — applied within a carefully designed human-in-the-loop architecture, can compress a four-hour manual search into a twelve-minute coordinated process. The platform provides a validated architectural and methodological foundation for future operational deployment in partnership with veterinary institutions and regulatory bodies across India.

\newpage
\begin{center}
    \textbf{\large APPENDICES}\\[0.6em]
    \textbf{\large APPENDIX A — SAMPLE CODING}
\end{center}
\vspace{1em}

\noindent The following code listings present representative excerpts from the K9Hope source repository (\url{https://github.com/Rvunveil/k9hope-canine-blood-donation}). Six listings are provided, covering the six functional areas of the platform: (A1) Firebase client configuration, (A2) authentication middleware and role-based routing, (A3) user login and role-differentiated database initialisation, (A4) donor eligibility and urgent request retrieval, (A5) Haversine geospatial distance computation and GMA query, and (A6) the ZICP training, calibration, and inference pipeline. All listings are drawn from production source files; minor inline comments have been added for report readability.

\vspace{1.5em}
\noindent\textbf{Listing A1: Firebase Client Configuration}
(\texttt{firebaseConfig.ts})
\vspace{0.4em}
\begin{lstlisting}[language=TypeScript, caption={Firebase client SDK initialisation with environment-variable-backed config. All sensitive keys are loaded from \texttt{.env.local} and never committed to the repository.}]
import { initializeApp, getApps, getApp } from "firebase/app";
import { getFirestore }  from "firebase/firestore";
import { getAuth }       from "firebase/auth";
import { getDatabase }   from "firebase/database";

const firebaseConfig = {
  apiKey:            process.env.NEXT_PUBLIC_FIREBASE_API_KEY,
  authDomain:        process.env.NEXT_PUBLIC_FIREBASE_AUTH_DOMAIN,
  projectId:         process.env.NEXT_PUBLIC_FIREBASE_PROJECT_ID,
  storageBucket:     process.env.NEXT_PUBLIC_FIREBASE_STORAGE_BUCKET,
  messagingSenderId: process.env.NEXT_PUBLIC_FIREBASE_MESSAGING_SENDER_ID,
  appId:             process.env.NEXT_PUBLIC_FIREBASE_APP_ID,
  databaseURL:       process.env.NEXT_PUBLIC_FIREBASE_DATABASE_URL,
};

// Prevent duplicate app initialisation during Next.js hot-reload
const app = getApps().length ? getApp() : initializeApp(firebaseConfig);

export const db   = getFirestore(app);
export const auth = getAuth(app);
export const rtdb = getDatabase(app);
\end{lstlisting}

\vspace{1.2em}
\noindent\textbf{Listing A2: Authentication Middleware and Role-Based Route Protection}
(\texttt{middleware.ts})
\vspace{0.4em}
\begin{lstlisting}[language=TypeScript, caption={Next.js Edge Middleware enforces cookie-based session presence before serving any protected route. Unauthenticated requests are redirected to \texttt{/login}. Role-specific path prefixes (\texttt{/app/h}, \texttt{/app/d}, \texttt{/app/p}, \texttt{/app/o}) map to Veterinary, Donor, Patient, and Organisation dashboards respectively.}]
import { NextResponse } from 'next/server';
import type { NextRequest } from 'next/server';

export function middleware(request: NextRequest) {
  const { pathname } = request.nextUrl;

  // Public route — allow unconditionally
  if (pathname === '/login') return NextResponse.next();

  const userIdCookie = request.cookies.get('userId');
  const roleCookie   = request.cookies.get('role');

  // Protected path prefixes: veterinary / donor / patient / organisation / onboarding
  const protectedRoutes = ['/app/h', '/app/d', '/app/p', '/app/o', '/onboarding'];
  const isProtected = protectedRoutes.some(r => pathname.startsWith(r));

  if (isProtected && (!userIdCookie || !roleCookie)) {
    // No valid session cookies — hard redirect to login
    return NextResponse.redirect(new URL('/login', request.url));
  }

  return NextResponse.next();
}

export const config = {
  matcher: ['/app/:path*', '/onboarding/:path*'],
};
\end{lstlisting}

\vspace{1.2em}
\noindent\textbf{Listing A3: User Login and Role-Differentiated Database Initialisation}
(\texttt{firebaseFunctions.ts} — \texttt{loginUserDatabase})
\vspace{0.4em}
\begin{lstlisting}[language=TypeScript, caption={The \texttt{loginUserDatabase} function resolves an incoming login credential (phone for Donor/Patient roles; email for Veterinary/Organisation roles) against the Firestore \texttt{users} collection. If the user is new, it creates parallel documents in the top-level \texttt{users} collection and the role-specific collection (\texttt{patients}, \texttt{donors}, \texttt{clinics}, or \texttt{organisations}). A custom 24-character random ID is generated to avoid exposing sequential integer primary keys.}]
// Generates a 24-character random alphanumeric user ID  (e.g. "aBc1-dEf2-gHi3-jKl4-mNo5-pQr6")
function generateUserId(): string {
  const chars = "ABCDEFGHIJKLMNOPQRSTUVWXYZabcdefghijklmnopqrstuvwxyz0123456789";
  return Array.from({ length: 6 }, () =>
    Array.from({ length: 4 }, () =>
      chars.charAt(Math.floor(Math.random() * chars.length))
    ).join("")
  ).join("-");
}

export async function loginUserDatabase(role: string, loginId: string) {
  const normalized = loginId.toLowerCase().trim();

  // --- DONOR role ---
  if (role === "donor") {
    const usersRef = collection(db, "users");
    const snap = await getDocs(query(usersRef, where("phone", "==", normalized)));

    if (!snap.empty) {
      const uid = snap.docs[0].id;
      // Ensure donor sub-document exists (idempotent upsert)
      await setDoc(doc(db, "donors", uid),
        { phone: normalized, onboarded: "no", createdAt: new Date() },
        { merge: true }
      );
      return uid;
    }

    // First-time donor: create user + donor documents atomically
    const uid = generateUserId();
    await setDoc(doc(db, "users", uid), {
      phone: normalized, roles: ["donor"],
      role: "donor", onboarded: "no",
      createdAt: new Date(), updatedAt: new Date()
    });
    await setDoc(doc(db, "donors", uid), {
      phone: normalized, onboarded: "no", createdAt: new Date()
    });
    return uid;
  }

  // --- VETERINARY role (email-based) ---
  if (role === "veterinary") {
    const usersRef = collection(db, "users");
    const snap = await getDocs(query(usersRef, where("email", "==", normalized)));
    if (!snap.empty) return snap.docs[0].id;

    const uid = generateUserId();
    await setDoc(doc(db, "users", uid), {
      email: normalized, role: "veterinary",
      onboarded: "no", createdAt: new Date()
    }, { merge: true });
    await setDoc(doc(db, "clinics", uid), {
      email: normalized, onboarded: "no",
      createdAt: new Date()
    }, { merge: true });
    return uid;
  }

  // Patient and Organisation follow the same pattern (omitted for brevity)
  throw new Error(`Unsupported role: ${role}`);
}
\end{lstlisting}

\vspace{1.2em}
\noindent\textbf{Listing A4: Urgent Request Retrieval with Blood-Type and City Matching}
(\texttt{firebaseFunctions.ts} — \texttt{getUrgentRequests})
\vspace{0.4em}
\begin{lstlisting}[language=TypeScript, caption={Retrieves onboarded patient records and performs client-side multi-constraint filtering: urgency tier (\texttt{immediate} or \texttt{within\_24\_hours}), DEA blood-type compatibility, and geographic proximity by city. This implements the presentation layer of the GMA result, surfacing the coordinator-reviewed shortlist on the Donor dashboard.}]
export async function getUrgentRequests(
  donorCity?: string,
  donorBloodType?: string,
  onlyUrgent: boolean = true
) {
  const patientsRef = collection(db, "patients");

  // Server-side filter: onboarded patients only (Firestore index: onboarded ASC)
  const snap = await getDocs(
    query(patientsRef, where("onboarded", "==", "yes"), limit(50))
  );

  // Cast to typed records and sort by submission time (descending) client-side
  let requests = snap.docs.map(d => ({ id: d.id, ...d.data() })) as any[];
  requests.sort((a, b) => {
    const tA = a.createdAt?.toDate?.() ?? new Date(0);
    const tB = b.createdAt?.toDate?.() ?? new Date(0);
    return tB.getTime() - tA.getTime();
  });

  // Urgency tier filter
  if (onlyUrgent) {
    requests = requests.filter(r =>
      r.p_urgencyRequirment === "immediate" ||
      r.p_urgencyRequirment === "within_24_hours"
    );
  }

  // Blood-type compatibility filter (exact DEA type match for MVP)
  if (donorBloodType) {
    requests = requests.filter(r => r.p_bloodgroup === donorBloodType);
  }

  // Prioritise same-city donors by stable re-sort
  if (donorCity) {
    requests.sort((a, b) => {
      const aLocal = a.p_city?.toLowerCase() === donorCity.toLowerCase() ? 0 : 1;
      const bLocal = b.p_city?.toLowerCase() === donorCity.toLowerCase() ? 0 : 1;
      return aLocal - bLocal;
    });
  }

  return requests.slice(0, onlyUrgent ? 6 : 20);
}
\end{lstlisting}

\vspace{1.2em}
\noindent\textbf{Listing A5: Haversine Geospatial Distance Computation and GMA Donor Shortlist}
(\texttt{lib/gma.ts} — illustrative implementation)
\vspace{0.4em}
\begin{lstlisting}[language=TypeScript, caption={The Haversine formula computes the great-circle distance between two GPS coordinates in kilometres. The GMA wraps it with a bounding-box pre-filter on Firestore \texttt{GeoPoint} fields to reduce the candidate set before exact distance computation, consistent with the algorithm specification in Chapter 3.}]
const R_EARTH_KM = 6371;

/**
 * Haversine great-circle distance (km) between two WGS-84 coordinates.
 * @param lat1  Latitude  of clinic  (degrees)
 * @param lon1  Longitude of clinic  (degrees)
 * @param lat2  Latitude  of donor   (degrees)
 * @param lon2  Longitude of donor   (degrees)
 */
export function haversineDistanceKm(
  lat1: number, lon1: number,
  lat2: number, lon2: number
): number {
  const toRad = (d: number) => (d * Math.PI) / 180;
  const dLat  = toRad(lat2 - lat1);
  const dLon  = toRad(lon2 - lon1);
  const a =
    Math.sin(dLat / 2) ** 2 +
    Math.cos(toRad(lat1)) * Math.cos(toRad(lat2)) *
    Math.sin(dLon / 2) ** 2;
  return R_EARTH_KM * 2 * Math.atan2(Math.sqrt(a), Math.sqrt(1 - a));
}

/**
 * Returns donors ranked by ascending Haversine distance from the clinic,
 * filtered to DAHD-eligible, cooldown-clear, DEA-compatible profiles
 * within the configured radius.
 */
export async function getGMAShortlist(
  clinicLat: number,
  clinicLon: number,
  requiredDeaType: string,
  radiusKm: number = 10
): Promise<Array<{ id: string; distanceKm: number; [k: string]: any }>> {

  // Bounding-box degrees approximation for Firestore pre-filter
  const dLat  = radiusKm / 111;
  const dLon  = radiusKm / (111 * Math.cos((clinicLat * Math.PI) / 180));

  const donorsRef = collection(db, "donors");
  const snap = await getDocs(
    query(
      donorsRef,
      where("eligibilityStatus", "==", "verified"),
      where("deaType", "==", requiredDeaType)
    )
  );

  const today  = new Date();
  const candidates = snap.docs
    .map(d => ({ id: d.id, ...d.data() } as any))
    // Exclude cooldown donors
    .filter(d => !d.cooldownUntil || d.cooldownUntil.toDate() < today)
    // Apply bounding-box pre-filter
    .filter(d => {
      const dLa = Math.abs(d.locationGeo.latitude  - clinicLat);
      const dLo = Math.abs(d.locationGeo.longitude - clinicLon);
      return dLa <= dLat && dLo <= dLon;
    })
    // Compute exact Haversine distance
    .map(d => ({
      ...d,
      distanceKm: haversineDistanceKm(
        clinicLat, clinicLon,
        d.locationGeo.latitude,
        d.locationGeo.longitude
      )
    }))
    // Final radius enforcement and ascending sort
    .filter(d => d.distanceKm <= radiusKm)
    .sort((a, b) => a.distanceKm - b.distanceKm);

  return candidates;
}
\end{lstlisting}

\vspace{1.2em}
\noindent\textbf{Listing A6: ZICP Training, Calibration, and Inference Pipeline}
(Python 3.10 — scheduled Firebase Cloud Function)
\vspace{0.4em}
\begin{lstlisting}[language=Python3, caption={Complete Zero-Inflated Conformal Prediction (ZICP) implementation as described in the IEEE ZICP paper. LightGBM models are trained for occurrence and log-normal size; non-conformity scores are computed on held-out calibration data; the empirical quantile threshold is used to construct distribution-free prediction intervals for the next 7 days per clinic per blood type.}]
import numpy as np
import pandas as pd
from lightgbm import LGBMClassifier, LGBMRegressor
from sklearn.preprocessing import LabelEncoder

# ── Non-conformity score (Equation 7 from IEEE ZICP paper) ──────────
def non_conformity_score(pi: float, mu: float, sigma: float, d: float) -> float:
    """
    Two-part score:
      D = 0  →  low score when model predicts high P(D=0)
      D > 0  →  low score when occurrence and log-size are jointly accurate
    """
    if d == 0:
        return -np.log(1 - pi + 1e-9)           # penalise high occurrence probability
    else:
        return -np.log(pi + 1e-9) + (np.log(d) - mu) ** 2 / (2 * sigma ** 2 + 1e-9)


class ZICP:
    """Zero-Inflated Conformal Prediction for sparse veterinary blood demand."""

    def __init__(self, alpha: float = 0.05):
        self.alpha    = alpha
        self.occ_mdl  = LGBMClassifier(n_estimators=200, learning_rate=0.05,
                                        max_depth=4, num_leaves=15,
                                        min_child_samples=10, subsample=0.8,
                                        colsample_bytree=0.8, random_state=42)
        self.size_mdl = LGBMRegressor(n_estimators=200, learning_rate=0.05,
                                       max_depth=4, num_leaves=15,
                                       min_child_samples=5, subsample=0.8,
                                       colsample_bytree=0.8, random_state=42)
        self.std_mdl  = LGBMRegressor(n_estimators=200, learning_rate=0.05,
                                       max_depth=4, num_leaves=15,
                                       min_child_samples=5, random_state=42)
        self.q_hat    = None   # calibrated conformal threshold

    # ── Training ────────────────────────────────────────────────────
    def fit(self, X_train: np.ndarray, y_train: np.ndarray):
        z_train = (y_train > 0).astype(int)
        self.occ_mdl.fit(X_train, z_train)

        pos_mask    = y_train > 0
        X_pos       = X_train[pos_mask]
        log_y_pos   = np.log(y_train[pos_mask])
        self.size_mdl.fit(X_pos, log_y_pos)

        mu_pos      = self.size_mdl.predict(X_pos)
        residuals   = np.abs(log_y_pos - mu_pos)
        self.std_mdl.fit(X_pos, residuals)

    # ── Calibration (Algorithm 1 from IEEE ZICP paper) ──────────────
    def calibrate(self, X_cal: np.ndarray, y_cal: np.ndarray):
        scores = []
        for xi, di in zip(X_cal, y_cal):
            pi    = self.occ_mdl.predict_proba(xi.reshape(1, -1))[0, 1]
            mu    = float(self.size_mdl.predict(xi.reshape(1, -1))[0])
            sigma = float(self.std_mdl.predict(xi.reshape(1, -1))[0])
            scores.append(non_conformity_score(pi, mu, sigma, di))

        n_cal     = len(scores)
        q_level   = np.ceil((n_cal + 1) * (1 - self.alpha)) / n_cal
        self.q_hat = np.quantile(scores, min(q_level, 1.0))

    # ── Inference (binary search for upper bound U) ──────────────────
    def predict_interval(self, x_new: np.ndarray) -> tuple[float, float]:
        pi    = self.occ_mdl.predict_proba(x_new.reshape(1, -1))[0, 1]
        mu    = float(self.size_mdl.predict(x_new.reshape(1, -1))[0])
        sigma = float(self.std_mdl.predict(x_new.reshape(1, -1))[0])

        U = 0
        for d_candidate in range(0, 21):         # demand bounded at 20 units
            if non_conformity_score(pi, mu, sigma, d_candidate) <= self.q_hat:
                U = d_candidate
        return (0.0, float(U))                   # L always 0 (non-negative demand)

    # ── Inventory order quantity (Equation 3 from IEEE ZICP paper) ───
    @staticmethod
    def order_quantity(upper: float, current_inventory: float) -> float:
        return max(0.0, upper - current_inventory)


# ── Scheduled Cloud Function entry point ────────────────────────────
def run_zicp_forecast(demand_records: pd.DataFrame) -> pd.DataFrame:
    """
    demand_records: DataFrame with columns
      [clinic_id, blood_type, date, demand, dow, month, rolling_7d_mean]

    Returns:
      DataFrame with columns [clinic_id, blood_type, forecast_date, lower, upper]
      for each clinic-blood-type combination, 7 days ahead.
    """
    le_clinic = LabelEncoder()
    le_btype  = LabelEncoder()
    demand_records["clinic_enc"] = le_clinic.fit_transform(demand_records["clinic_id"])
    demand_records["btype_enc"]  = le_btype.fit_transform(demand_records["blood_type"])

    feature_cols = ["dow", "month", "clinic_enc", "btype_enc", "rolling_7d_mean"]
    demand_records.sort_values("date", inplace=True)

    n = len(demand_records)
    i_train = int(n * 0.50)
    i_cal   = int(n * 0.75)

    X        = demand_records[feature_cols].values
    y        = demand_records["demand"].values
    X_train, y_train = X[:i_train],       y[:i_train]
    X_cal,   y_cal   = X[i_train:i_cal],  y[i_train:i_cal]

    model = ZICP(alpha=0.05)
    model.fit(X_train, y_train)
    model.calibrate(X_cal, y_cal)

    results = []
    last_row     = demand_records.iloc[-1]
    last_date    = pd.to_datetime(last_row["date"])
    for clinic_id in demand_records["clinic_id"].unique():
        for blood_type in demand_records["blood_type"].unique():
            for k in range(1, 8):
                forecast_date = last_date + pd.Timedelta(days=k)
                x_new = np.array([
                    forecast_date.dayofweek,
                    forecast_date.month,
                    le_clinic.transform([clinic_id])[0],
                    le_btype.transform([blood_type])[0],
                    demand_records[
                        (demand_records["clinic_id"] == clinic_id) &
                        (demand_records["blood_type"] == blood_type)
                    ]["demand"].iloc[-7:].mean()
                ], dtype=float)
                lower, upper = model.predict_interval(x_new)
                results.append({
                    "clinic_id": clinic_id, "blood_type": blood_type,
                    "forecast_date": forecast_date.date(),
                    "lower": lower, "upper": upper
                })

    return pd.DataFrame(results)
\end{lstlisting}

\newpage
\begin{center}
    \textbf{\large REFERENCES}
\end{center}
\vspace{1em}

\begin{thebibliography}{99}

\bibitem{tocci2009}
L.~J. Tocci,
``Transfusion medicine in small animal practice,''
\textit{Veterinary Clinics of North America: Small Animal Practice},
vol.~40, no.~3, pp.~485--494, 2009.

\bibitem{brooks2020}
M.~B. Brooks et al.,
``2020 AAHA Canine Blood Component Transfusion Guidelines,''
\textit{Journal of the American Animal Hospital Association},
vol.~56, no.~1, pp.~1--21, 2020.

\bibitem{davidow2013}
E.~B. Davidow et al.,
``Association of Veterinary Hematology and Transfusion Medicine (AVHTM)
Transfusion Reaction Small Animal Consensus Statement (TRACS),''
\textit{Transfusion Medicine and Hemotherapy},
vol.~40, no.~5, pp.~322--334, 2013.

\bibitem{sharma2014}
A.~Sharma and K.~Raj,
``Dog erythrocyte antigen 1.1 significance in canine blood transfusion,''
\textit{Veterinary World},
vol.~7, no.~7, pp.~485--488, 2014.

\bibitem{mcmichael2009}
M.~A. McMichael, S.~A. Smith, and A.~Galligan,
``Effect of leukoreduction on transfusion-induced immunomodulation in dogs,''
\textit{Journal of Veterinary Emergency and Critical Care},
vol.~19, no.~1, pp.~65--70, 2009.

\bibitem{lanevschi2001}
A.~Lanevschi and K.~J. Wardrop,
``Principles of transfusion medicine in small animals,''
\textit{Canadian Veterinary Journal},
vol.~42, no.~6, pp.~447--454, 2001.

\bibitem{kumar2022}
R.~Kumar, P.~Singh, and A.~Sharma,
``Seasonal variations in canine vector-borne diseases in metropolitan India,''
\textit{Indian Journal of Veterinary Sciences},
vol.~15, no.~4, pp.~345--358, 2022.

\bibitem{spada2016}
E.~Spada et al.,
``Seroprevalence of canine vector-borne pathogens in northern Italy,''
\textit{Parasites \& Vectors},
vol.~9, no.~1, pp.~1--10, 2016.

\bibitem{dini2025}
R.~Dini et al.,
``Insights into the canine blood donor experience: A multicenter study
on physiological and behavioral changes,''
\textit{Veterinary Sciences},
vol.~12, no.~9, p.~876, 2025,
doi: \texttt{10.3390/vetsci12090876}.

\bibitem{dahd2025}
Department of Animal Husbandry and Dairying, Ministry of Fisheries, Animal
Husbandry and Dairying, Government of India,
``Guidelines/SOP for Blood Transfusion \& Blood Bank for Animals in India
(Draft for Stakeholder Comment),''
23 July 2025. [Online]. Available:
\url{https://avantiscdn.prod.storage.blob.core.windows.net/legalupdatedocs/45117DAHD-issued-the-Draft-Guidelines-SOP-for-Blood-Transfusion-Blood-Bank-for-Animals-in-India-July312025.pdf}

\bibitem{pib2025}
Press Information Bureau, Ministry of Information and Broadcasting,
Government of India,
``Fact Sheet: Guidelines SOP for Animal Blood Transfusion and Blood
Banks in India,''
26 August 2025. [Online]. Available:
\url{https://static.pib.gov.in/WriteReadData/specificdocs/documents/2025/aug/doc2025825619801.pdf}

\bibitem{zahiri2018}
B.~Zahiri and M.~S. Pishvaee,
``Blood supply chain network design considering blood group compatibility
under uncertainty,''
\textit{International Journal of Production Research},
vol.~56, no.~7, pp.~2486--2507, 2018.

\bibitem{butt2018}
M.~Butt, A.~Entezari, A.~Foad, and S.~Caragher,
``Blood transfusion demand and allocation optimization in emergency
situations,''
\textit{International Journal of Healthcare Management},
vol.~11, no.~3, pp.~201--215, 2018.

\bibitem{dehghani2021}
M.~Dehghani, B.~Abbasi, and F.~Oliveira,
``Proactive transshipment in the blood supply chain: A stochastic
programming approach,''
\textit{Omega},
vol.~98, article 102112, 2021,
doi: \texttt{10.1016/j.omega.2019.102112}.

\bibitem{khakshouri2024}
S.~Khakshouri Fariman, K.~Danesh, M.~Pourtalebiyan, Z.~Fakhri,
A.~Motallebi, and A.~Fozooni,
``A robust optimization model for multi-objective blood supply chain
network considering scenario analysis under uncertainty,''
\textit{Scientific Reports},
vol.~14, p.~9452, Apr.\ 2024,
doi: \texttt{10.1038/s41598-024-57521-0}.

\bibitem{croston1972}
J.~D. Croston,
``Forecasting and stock control for intermittent demands,''
\textit{Operational Research Quarterly},
vol.~23, no.~3, pp.~289--303, 1972.

\bibitem{liu2020}
P.~Liu,
``Intermittent demand forecasting for medical consumables with short life
cycle using a dynamic neural network during the COVID-19 epidemic,''
\textit{Health Informatics Journal},
vol.~26, no.~4, pp.~3106--3122, Dec.\ 2020,
doi: \texttt{10.1177/1460458220954730}.

\bibitem{dinamarca2025}
M.~A. Dinamarca, F.~Rojas, C.~Ibacache-Quiroga, and K.~Gonzalez-Pizarro,
``Modeling time series with SARIMAX and skew-normal and zero-inflated
skew-normal errors,''
\textit{Mathematics},
vol.~13, no.~11, p.~1892, 2025,
doi: \texttt{10.3390/math13111892}.

\bibitem{mutimbanyoka2025}
N.~H. Mutimbanyoka and A.~Ndlovu,
``Predictive modelling of blood demand,''
\textit{International Journal of Computer Science and Mobile Computing},
vol.~14, no.~6, pp.~26--33, Jun.\ 2025,
doi: \texttt{10.47760/ijcsmc.2025.v14i06.004}.

\bibitem{vovk2005}
V.~Vovk, A.~Gammerman, and G.~Shafer,
\textit{Algorithmic Learning in a Random World}.
New York: Springer, 2005.

\bibitem{angelopoulos2021}
A.~N. Angelopoulos and S.~Bates,
``A gentle introduction to conformal prediction and distribution-free
uncertainty quantification,''
\textit{arXiv preprint arXiv:2107.07511}, 2021.

\bibitem{romano2019}
Y.~Romano, E.~Patterson, and E.~Candes,
``Conformalized quantile regression,''
\textit{Advances in Neural Information Processing Systems},
vol.~32, pp.~3543--3553, 2019.

\bibitem{stankeviciute2021}
K.~Stankeviciute, A.~M. Alaa, and M.~van der Schaar,
``Conformal time-series forecasting,''
\textit{Advances in Neural Information Processing Systems},
vol.~34, pp.~6216--6228, 2021.

\bibitem{xu2023}
C.~Xu and Y.~Xie,
``Conformal prediction for time series,''
\textit{IEEE Transactions on Pattern Analysis and Machine Intelligence},
vol.~45, no.~10, pp.~11575--11587, Oct.\ 2023,
doi: \texttt{10.1109/TPAMI.2023.3272339}.

\bibitem{zaffran2022}
M.~Zaffran, A.~Dieuleveut, J.~Josse, and Y.~Goude,
``Adaptive conformal predictions for time series,''
in \textit{Proc.\ 39th Int.\ Conf.\ Machine Learning},
PMLR, vol.~162, 2022, pp.~25834--25866.

\bibitem{lei2018}
J.~Lei, M.~G'Sell, A.~Rinaldo, R.~J. Tibshirani, and L.~Wasserman,
``Distribution-free predictive inference for regression,''
\textit{Journal of the American Statistical Association},
vol.~113, no.~523, pp.~1094--1111, 2018.

\bibitem{vazquez2022}
J.~Vazquez and J.~C. Facelli,
``Conformal prediction in clinical medical sciences,''
\textit{Journal of Healthcare Informatics Research},
vol.~6, no.~3, pp.~241--252, Sep.\ 2022,
doi: \texttt{10.1007/s41666-021-00113-8}.

\bibitem{yeh2025}
C.-H. Yeh, C.-M. Hung, H.-H. Chung, and W.-J. Su,
``Prediction of blood donor return behaviour using machine learning,''
\textit{Vox Sanguinis}, 2025,
doi: \texttt{10.1111/vox.70139},
PMID: 41115732, Online ahead of print.

\bibitem{chawla2002}
N.~V. Chawla, K.~W. Bowyer, L.~O. Hall, and W.~P. Kegelmeyer,
``SMOTE: Synthetic minority over-sampling technique,''
\textit{Journal of Artificial Intelligence Research},
vol.~16, pp.~321--357, 2002.

\bibitem{cornell2023}
Cornell University College of Veterinary Medicine,
``Transfusion Guidelines,''
Animal Health Diagnostic Center, Comparative Coagulation Section,
last updated February 2023. [Online]. Available:
\url{https://www.vet.cornell.edu/animal-health-diagnostic-center/laboratories/comparative-coagulation/clinical-topics/transfusion-guidelines}

\bibitem{kohn2006}
B.~Kohn, C.~Weingart, V.~Eckmann, M.~Ottenjann, and W.~Leibold,
``Primary immune-mediated hemolytic anemia in 19 dogs: diagnosis, therapy,
and outcome,''
\textit{Journal of Veterinary Internal Medicine},
vol.~20, no.~3, pp.~159--166, 2006.

\bibitem{prittie2010}
J.~Prittie,
``Triggers for use, testing, and transfusion of blood and blood products,''
\textit{Veterinary Clinics of North America: Small Animal Practice},
vol.~40, no.~3, pp.~465--476, 2010.

\bibitem{cdfa2025}
California Department of Food and Agriculture,
``Animal Blood Banking Document 2.0 (DRAFT),''
Animal Health and Food Safety Services, Commercial Animal Blood Banks
Program, 3 April 2025. [Online]. Available:
\url{https://www.cdfa.ca.gov/AHFSS/cabb/docs/animal_blood_banking_guidance_resource.pdf}

\bibitem{pandithurai2026a}
O. Pandithurai, A. Akash, J. Praveenkumar, T. Vikram, D. Prem Kumar, and P. Ram Kishore, 
``Decision Support for Canine Blood Services: A Platform Architecture for Sparse Veterinary Networks,'' 
\textit{Rajalakshmi Institute of Technology, Chennai}, IEEE Manuscript, Feb. 2026.

\bibitem{pandithurai2026b}
O. Pandithurai, A. Akash, J. Praveenkumar, T. Vikram, D. Prem Kumar, and P. Ram Kishore, 
``Conformal Prediction for Zero-Inflated Blood Demand: Uncertainty-Aware Forecasting in Sparse Veterinary Networks,'' 
\textit{Rajalakshmi Institute of Technology, Chennai}, IEEE Manuscript, 2026.

\bibitem{scconline2025}
SCC Online, ``First-ever Guidelines for Animal Blood Transfusion Released,'' 
\textit{SCC Times}, 26 August 2025. [Online]. Available: 
\url{https://www.scconline.com/blog/post/2025/08/26/first-ever-animal-blood-transfusion-guidelines-india-2025}

\bibitem{pcaa1960}
Department of Animal Husbandry and Dairying, ``Prevention of Cruelty to Animals Act, 1960 (as amended),'' 
Ministry of Fisheries, Animal Husbandry and Dairying, Government of India.

\bibitem{ivca1984}
Indian Veterinary Council Act, 1984, 
Ministry of Fisheries, Animal Husbandry and Dairying, Government of India.

\end{thebibliography}

\end{document}
